%% Q14 — Toward Fermionic Matter from Projective Dirac Admissibility
%% Cosmochrony programme, Q-series paper 14
%% J. Beau, 2026

\documentclass[11pt,a4paper]{article}

\usepackage[T1]{fontenc}
\usepackage[utf8]{inputenc}
\usepackage{amsmath,amssymb,amsthm}
\usepackage[colorlinks=true,linkcolor=blue,citecolor=blue,urlcolor=blue]{hyperref}
\hypersetup{
  pdftitle={Toward Fermionic Matter from Projective Dirac Admissibility: The Electroweak Spinor Algebra,
    and Chirality Conditional on a Lorentzian Spin Structure},
  pdfauthor={J\'er\^ome Beau},
  pdfsubject={Fermions read from the admissible Weil module: the electroweak spinor algebra is proved;
    chiral selection and hypercharge constraints are conditional on a Lorentzian spin structure},
  pdfkeywords={admissible spinor bundle, metaplectic lift, non-injective projection,
    Weil representation, Lorentzian spin structure, projected Dirac operator, chiral selection, hypercharge rigidity,
    generation multiplicity, spectral anomaly cancellation, Cosmochrony}
}
\usepackage{natbib}
\usepackage{geometry}
\geometry{top=2.5cm,bottom=2.5cm,left=2.8cm,right=2.8cm}
\usepackage{microtype}
\usepackage{mathtools}
\usepackage{doi}

%% Theorem environments
\newtheorem{theorem}{Theorem}[section]
\newtheorem{lemma}[theorem]{Lemma}
\newtheorem{proposition}[theorem]{Proposition}
\newtheorem{corollary}[theorem]{Corollary}
\theoremstyle{definition}
\newtheorem{definition}[theorem]{Definition}
\newtheorem{hypothesis}[theorem]{Hypothesis}
\theoremstyle{remark}
\newtheorem{remark}[theorem]{Remark}

%% Math operators
\DeclareMathOperator{\Sym}{Sym}
\DeclareMathOperator{\tr}{tr}
\DeclareMathOperator{\Tr}{Tr}
\DeclareMathOperator{\End}{End}
\DeclareMathOperator{\Hom}{Hom}
\DeclareMathOperator{\ad}{ad}
\DeclareMathOperator{\Mult}{Mult}
\newcommand{\Mp}{\operatorname{Mp}}
\newcommand{\Spin}{\operatorname{Spin}}
\newcommand{\SL}{\operatorname{SL}}
\newcommand{\SU}{\operatorname{SU}}
\newcommand{\Rep}{\operatorname{Rep}}

\title{Toward Fermionic Matter from Projective Dirac Admissibility:\\
The Electroweak Spinor Algebra, and Chirality\\
Conditional on a Lorentzian Spin Structure}

\author{J\'er\^ome Beau\\[4pt]
\small Independent Researcher, France\\
\small \texttt{jerome.beau@cosmochrony.org}}

\date{2026}

\begin{document}

  \maketitle

  \begin{abstract}
    The gauge--gravity synthesis of the Cosmochrony programme reads gravity and Yang--Mills dynamics,
    conditionally, as the $a_2$ and $a_4$ Seeley--DeWitt responses of one admissible spectral functional.
    This paper asks what the fermionic sector requires when fermions are read from the Weil module of the
    admissible fibre on a supplied Heisenberg carrier, which every result below inherits.

    The algebraic part is proved.
    The complexified metaplectic algebra is
    $\mathfrak{mp}(2,\mathbb{R})_{\mathbb{C}} \simeq \mathfrak{sl}_2(\mathbb{C})$; on the abstract doublet
    its symmetric square is the adjoint and its exterior square the trivial line, and a Hermitian form
    selects the compact real form $\mathfrak{su}(2)$.
    Given the supplied Born--Infeld parity datum, the internal parity $HK$ satisfies $(HK)^2 = -1$ without
    any metric.
    The physical readings require two distinct hypotheses, stated in the paper and supplied by no source.
    [H-Spin] identifies the $\SL(2,\mathbb{C})$ acting on the doublet with the spin group of a
    four-dimensional Lorentzian co-metric on the effective base; the geometric branch supplies such a
    co-metric only conditionally, and nothing relates the metaplectic action to its frames.
    [H-Weak] supplies a distinct rank-two weak factor: by Schur's lemma a weak action commuting with
    Lorentz transformations cannot live on the same copy of the doublet, so the symmetric square of the
    spinor factor is a Lorentz sector and its exterior square carries no hypercharge.
    Under [H-Spin], the projected Dirac operator contains a canonical zero-order endomorphism $E_\Pi$, the
    spectral residue of non-injective projection; the spinorial lift of the parity is unique and reverses
    chirality, and $E_\Pi$ is left-admissible on the orientation-compatible branch.
    The chiral selection of the weak interaction, the $V{-}A$ structure, and the anomaly-cancellation
    constraints on the hypercharge weights require [H-Weak] as well, and left-admissibility does not
    select the chiral assignment.
    Both hypotheses are missing identifications, not refutations.
    The supplied rank-three selection rule $\sigma_c(n_3) = 3$ (O23) admits a spinorial multiplicity
    reading, giving conditionally a gauge-singlet three-generation factor
    $\mathbb{C}^3_{\mathrm{gen}} \subset \ker(\operatorname{ad}_{\SU(2)} \oplus Y)$.
    The quark sector uses a supplied colour module; no $\mathrm{SU}(3)$ is derived.
    Finally, in a metaplectic step model the ordered step generator $\log g$ has a non-zero $J_3$
    component in the sector that can lift the static degeneracy of the generation factor; its
    identification with an emergent ordering derivative is not supplied, and the amplitude is open.
    Interpretive outlook (a reading, not a result): chirality is where the internal algebra of the
    admissible fibre must meet spacetime geometry and the weak interaction at once.
    The paper locates that meeting point in two explicit hypotheses, which turns the question of why
    weak interactions are left-handed into the question of how the metaplectic symmetry of the fibre is
    joined to Lorentz frames and to a separate internal doublet.
  \end{abstract}

  \medskip
  \noindent\textbf{Keywords:}
  admissible spinor bundle;
  metaplectic lift;
  non-injective projection;
  Weil representation;
  Lorentzian spin structure;
  projected Dirac operator;
  chiral selection;
  hypercharge rigidity;
  generation multiplicity;
  spectral anomaly cancellation;
  Cosmochrony

  \newpage
  \tableofcontents

  \newpage
  \section{Introduction}
    \label{sec:introduction}

    \subsection{Position of the problem}
      \label{ssec:position}

      The preceding papers of the programme propose a bosonic spectral stratification
      of admissible non-injective projection.
      In the Gravity trilogy~\cite{Beau2026h,Beau2026q12,Beau2026q13}, the horizontal $a_2$
      response gives an Einstein--Hilbert contribution in an explicit cutoff scheme, with the Newton
      constant fixed by a matching condition, the vertical $a_4$ response gives a Yang--Mills sector
      for a supplied gauge group, and the $a_6$ coefficient gives a list of candidate mixed
      gauge--gravity invariants, none of which is singled out.
      Reading the Seeley--DeWitt order as the organising principle of emergent interactions is an
      interpretive outlook of that synthesis, not a theorem~\cite{Beau2026q13}.
      What lies outside this bosonic synthesis is the fermionic sector: the structural
      origin of spinors, chirality, the distribution of hypercharge, and the generation
      multiplicity.

      The purpose of the present paper is to show that fermions need not be introduced
      as external matter fields added to the admissible bundle.
      They are read from the Weil module carried by the admissible fibre on a supplied Heisenberg
      carrier, through the complexification of its metaplectic Lie algebra.
      That complexification is algebra and needs no metric.
      Reading its integrated group as the spin group of an emergent spacetime is a separate step: it
      requires a four-dimensional Lorentzian co-metric and a spin structure identified with the
      metaplectic action, which this paper states once as Hypothesis~\ref{hyp:spin-solder} ([H-Spin]).
      The geometric branch of the programme supplies the co-metric only
      conditionally~\cite{Beau2026q5b,BeauQ8,Beau2026q11}, and no source supplies the identification.
      Reading the same algebra as the weak doublet is a third step, and it cannot use the same copy of
      the doublet: a weak action commuting with Lorentz transformations needs a distinct rank-two
      factor, stated as Hypothesis~\ref{hyp:weak-factor} ([H-Weak]) and not constructed here.

    \subsection{Main structural chain}
      \label{ssec:chain}

      The structural chain studied in this paper is
      \begin{equation}
        \label{eq:main-chain}
        \Pi
        \;\Rightarrow\;
        F_n \simeq V_\rho
        \;\Rightarrow\;
        \Mp(2,\mathbb{R})
        \;\Rightarrow\;
        \mathfrak{mp}(2,\mathbb{R})_{\mathbb{C}} \simeq \mathfrak{sl}_2(\mathbb{C})
        \;\overset{\text{[H-Spin]}}{\leadsto}\;
        \Spin(3,1) \simeq \SL(2,\mathbb{C})
        \;\Rightarrow\;
        S_\Pi .
      \end{equation}
      The identification $F_n \simeq V_\rho$ is the Weil-module characterisation of the
      admissible fibre on a \emph{supplied} Heisenberg carrier.
      It is not derived from the axioms: Foundation~\cite{Beau2026m} and
      HeisenbergStructure~\cite{Beau2026HeisenbergStructure} establish that A1--A3 force
      irreducibility and non-commutation, but neither makes the commutator central nor selects a
      finite Heisenberg group, an explicit $\mathfrak{S}_3$ countermodel satisfying the extracted
      carrier contract while violating that conclusion.
      Stone--von Neumann is a conditional endpoint, supplying uniqueness at a fixed non-trivial
      central character once that carrier is given.
      The passage to $\Mp(2,\mathbb{R})$ is the metaplectic lift of the symplectic structure
      canonically carried by $V_\rho$.
      The complexification
      $\mathfrak{mp}(2,\mathbb{R})_{\mathbb{C}} \simeq \mathfrak{sl}_2(\mathbb{C})$ is a statement about
      Lie algebras and holds without any metric; its simply connected integration is
      $\SL(2,\mathbb{C})$.
      The arrow marked [H-Spin] is not an implication: identifying this $\SL(2,\mathbb{C})$ with the
      spin group $\Spin(3,1)$ of a Lorentzian co-metric on the effective base, and hence defining the
      admissible spinor bundle $S_\Pi$ as a spinor bundle of spacetime, is the content of
      Hypothesis~\ref{hyp:spin-solder}.

      On the spinor factor, $\operatorname{Sym}^2(S_L)$ is a Lorentz sector and $\wedge^2(S_L)$ is a
      trivial line; neither is electroweak data.
      Under [H-Weak], the same tensor functors applied to the distinct weak factor
      $E_{\mathrm{weak}}$ define the electroweak bundle data:
      \begin{equation}
        \label{eq:tensor-functors}
        \operatorname{Sym}^2(E_{\mathrm{weak}}) \;\Rightarrow\; \operatorname{ad}_{\mathbb{C}}(\mathfrak{su}(2)_L),
        \qquad
        \wedge^2(E_{\mathrm{weak}}) \;\Rightarrow\; L_Y .
      \end{equation}
      The symmetric sector gives the complex weak adjoint; after selecting the compact real
      form, it defines the $\SU(2)_L$ gauge sector.
      The determinant line $L_Y = \wedge^2(E_{\mathrm{weak}})$ supports the $U(1)_Y$ factor, non-trivially
      only because the structure group of $E_{\mathrm{weak}}$ is $U(2)$, a choice made in [H-Weak].
      The distribution of hypercharge weights $Y_R$ is not assumed: it is constrained
      by requiring that the projected spectral functional
      $S_\Pi[g, A, \mathcal{D}_{\Pi,g,A}]$ be invariant under $U(1)_Y$ gauge
      transformations.

    \subsection{Three modular results}
      \label{ssec:results}

      The paper is organised around three independent but compatible results.

      \paragraph{Theorem A (Spinorial electroweak bundle).}
        The complexified metaplectic algebra of the admissible Weil module is
        $\mathfrak{sl}_2(\mathbb{C})$, and on the abstract doublet the symmetric and determinant
        tensor sectors are the adjoint and the trivial line; a Hermitian form selects the compact real
        form $\mathfrak{su}(2)$.
        This part is algebraic and proved.
        Two physical readings are conditional and distinct: the reading of $S_\Pi$ as a $\Spin(3,1)$
        spinor bundle of the effective spacetime needs Hypothesis~\ref{hyp:spin-solder}, and the
        $\SU(2)_L \times U(1)_Y$ electroweak data~\eqref{eq:tensor-functors} need the distinct weak factor of
        Hypothesis~\ref{hyp:weak-factor}.

      \paragraph{Theorem B (Chiral selection and hypercharge rigidity).}
        The projected Dirac operator satisfies the Lichnerowicz-type identity
      \begin{equation}
        \label{eq:lichnerowicz}
        \mathcal{D}_{\Pi,g,A}^2
        = -(\nabla^{S,A})^2
        + \frac{R}{4}
        + \frac{1}{2}\gamma^\mu\gamma^\nu F^a_{\mu\nu} T^a_R
        + E_\Pi ,
        \end{equation}
        where $E_\Pi$ is the internal spectral residue of non-injective projection in the
        spinorial sector.
        Under the supplied Born--Infeld datum, $E_\Pi$ is left-admissible, $P_R E_\Pi P_R = 0$: a
        statement about the projective residue and the Lorentz chirality.
        The chiral selection of the weak coupling, the $V{-}A$ structure, requires in addition the
        left-handed doublet assignment of Hypothesis~\ref{hyp:weak-factor}, and no argument in this
        paper derives that assignment from left-admissibility.
        The same inputs make $U(1)_Y$ invariance impose anomaly-cancellation constraints on the
        hypercharge assignments $Y_R$.
        This result is spectral and variational; it depends on the gauge sector
        of~\cite{Beau2026q6a,Beau2026q12} and on the gauge--gravity synthesis
        of~\cite{Beau2026q13} for the coupling of the Dirac sector to the joint spectral
        geometry.
        It is conditional on Hypothesis~\ref{hyp:spin-solder}: the spin connection, the Clifford
        symbol and the chirality operator $\gamma_5$ all presuppose the Lorentzian spin structure; its
        gauge part is conditional on Hypothesis~\ref{hyp:weak-factor}.

      \paragraph{Theorem C (Generation multiplicity).}
        The admissible saturation invariant $\sigma_c(n_3) = 3$~\cite{Beau2026a27}, a
        supplied selection rule (O23 Theorem~3.1 proves the three-dimensionality of the
        neutral sector of a supplied spinor carrier; the carrier selection and the
        identification of $\Sigma_c$ with that dimension are open),
        admits a spinorial multiplicity reading.
        It yields a gauge-singlet generation space
      \begin{equation}
        \label{eq:generation-space}
        \mathbb{C}^3_{\mathrm{gen}}
        \;\subset\;
        \ker\!\left(\operatorname{ad}_{\SU(2)} \oplus Y\right),
        \end{equation}
        which is not identified with $\operatorname{Sym}^2(S_\Pi)$.
        The generation factor is a multiplicity space, distinct from the adjoint triplet and from any
        spatial triplet (Section~\ref{sec:generations}).
        This result is combinatorial and representational; it is conditional on the supplied
        rank-three selection rule of~\cite{Beau2026a27} (Section~4.2; cf.\ Theorem~3.1).

  \subsection{Status and dependencies}
    \label{ssec:status}

    The full matter content targeted by the construction is
    \begin{equation}
      \label{eq:matter-content}
      S_\Pi^{\mathrm{matter}}
      = \bigl(S_\Pi \oplus (S_\Pi \otimes V_{\mathrm{color}})\bigr)
      \otimes \mathbb{C}^3_{\mathrm{gen}},
    \end{equation}
    where $V_{\mathrm{color}} \in \Rep(\SU(3))$.
    The following status distinctions apply throughout the paper.
    The electroweak spinorial sector is structural once Theorem~A is established, and Theorem~A
    itself is conditional on the supplied Heisenberg carrier that fixes $F_n \simeq V_\rho$
    (Subsection~\ref{ssec:weil-metaplectic}); that condition is upstream of, and independent of,
    the colour module below.
    The colour-coupled quark bundle $S_\Pi \otimes V_{\mathrm{color}}$ is \emph{conditional} on a
    supplied colour module: O31~\cite{Beau2026a35} is a withdrawal notice, and
    $[H\text{-color}]_{\mathrm{pointwise}}$ is among the results it withdraws
    (Section~\ref{ssec:colour-status}).
    The hypercharge constraints of Theorem~B and the generation multiplicity of
    Theorem~C are derived below from their stated hypotheses; they are not assumed as inputs.

    Hypotheses~\ref{hyp:spin-solder} ([H-Spin]) and~\ref{hyp:weak-factor} ([H-Weak]) separate the paper
    into three layers.
    The algebraic layer uses neither: the complexification
    $\mathfrak{mp}(2,\mathbb{R})_{\mathbb{C}} \simeq \mathfrak{sl}_2(\mathbb{C})$, the fibrewise tensor
    identities of Theorem~\ref{thm:spinorial-electroweak-bundle}(a), the grading identity of
    Lemma~\ref{lem:bi-grading-anticommute} and the quaternionic square $(HK)^2 = -1$ of the internal
    parity given the supplied Born--Infeld datum, and the linear algebra of
    Propositions~\ref{prop:dg-static}, \ref{prop:dg-odd} and~\ref{prop:dg-bridge} as statements about
    operators on $\mathbb{C}^3_{\mathrm{gen}}$.
    The Lorentz-covariant layer uses [H-Spin]: the spinor bundle as a bundle of spacetime spinors, the
    projected Dirac operator and its Lichnerowicz form, the chirality operator $\gamma_5$, the
    covariant structure of Lemma~\ref{lem:covariant-epsilon}, the uniqueness and chirality reversal
    of the spinorial BI lift (Theorem~\ref{thm:spinorial-bi-lift}), left-admissibility, and the reading
    of the operators of Section~\ref{sec:dyn-lift} as restrictions of $E_\Pi^2$.
    The electroweak layer uses both: the $\SU(2)_L \times U(1)_Y$ bundle data, the $V{-}A$ structure of
    weak couplings, and the anomaly constraints on the hypercharge weights.
    Both hypotheses are missing identifications, not refutations: no result excludes them, and each
    layer holds as stated once its hypotheses are supplied.
    Relative to them, the $V{-}A$ structure of Theorem~B depends on the supplied Heisenberg carrier,
    the axiomatic chain A1--A4~\cite{Beau2026m}, the Born--Infeld datum supplied as hypotheses
    (H1)--(H2) of O30~\cite{Beau2026a34}, and the chiral assignment of
    Hypothesis~\ref{hyp:weak-factor}; it does not depend on the colour module.
    The hypercharge rigidity of Theorem~B is not: the anomaly
    constraints~\eqref{eq:anomaly-linear}--\eqref{eq:anomaly-mixed} are weighted by $\dim(R)$, and
    the cubic constraint~\eqref{eq:anomaly-cubic} is the one that carries the colour multiplicity
    --- on the Standard Model pattern it vanishes precisely at three, by direct substitution
    in~\eqref{eq:anomaly-cubic}.
    This paper does not analyse the constraint system at general multiplicity and does not claim
    that the constraints force three; the multiplicity is a supplied input
    (Section~\ref{ssec:colour-status}).

    This paper depends simultaneously on four layers of the programme: the admissible
    Weil fibre structure (Foundation, HeisenbergStructure, O-series); the effective co-metric
    of the geometric branch (Q5b, Q8, Q11), which enters only through
    Hypothesis~\ref{hyp:spin-solder}; the gauge connection (Q6a, Q12); and the
    gauge--gravity synthesis (Q13, Gravity).
    It is accordingly the first paper of the Q-series whose construction depends on
    the simultaneous use of all four layers.

  \subsection{Organisation}
    \label{ssec:organisation}

    Section~\ref{sec:metaplectic} constructs the metaplectic lift and the admissible
    spinor bundle $S_\Pi$, proving Theorem~A and stating Hypothesis~\ref{hyp:spin-solder}.
    Section~\ref{sec:dirac} defines the projected Dirac operator, identifies $E_\Pi$
    as the spectral residue of non-injective projection in the spinorial sector, and
    proves the chiral selection result of Theorem~B.
    Section~\ref{sec:anomaly} derives the hypercharge rigidity from spectral anomaly
    cancellation.
    Section~\ref{sec:generations} establishes the spinorial multiplicity reading of
    $\sigma_c(n_3) = 3$ (a supplied selection rule) and proves Theorem~C.
    Section~\ref{sec:dyn-lift} studies the first structural constraint on the spectrum of
    $E_\Pi^2$ on the generation factor: a static $J_\Pi$-real, weight-preserving restriction
    cannot split the outer pair, whereas the per-step generator of an ordered metaplectic step
    carries a non-zero $J_3$ component capable of lifting this degeneracy.
    Section~\ref{sec:colour} discusses the colour-coupled quark sector and the status of the
    supplied colour module.
    Section~\ref{sec:conclusion} summarises the results and identifies the remaining
    open directions.

    \newpage
  \section{Metaplectic Lift and the Admissible Spinor Bundle}
  \label{sec:metaplectic}

  \subsection{From the Weil module to the metaplectic lift}
    \label{ssec:weil-metaplectic}

    The admissible fibre $F_n$ is identified with the Weil module $V_\rho$ of
    $\mathrm{Heis}_3(\mathbb{Z}/q\mathbb{Z})$ at a fixed non-trivial central character, on a
    Heisenberg carrier that is supplied rather than selected by the
    axioms~\cite{Beau2026m,Beau2026HeisenbergStructure}.
    This identification is not merely a representation-theoretic convenience.
    The Weil module carries the canonical symplectic structure associated with the
    Heisenberg commutator.
    Consequently, the automorphisms preserving the admissible commutator act through
    the metaplectic lift:
    \[
      F_n \simeq V_\rho
      \quad\Rightarrow\quad
      \Mp(2,\mathbb{R}).
    \]
    The passage to $\Mp(2,\mathbb{R})$ is therefore structural: it is the double cover
    naturally associated with the symplectic action underlying the Weil representation.

  \subsection{Complexification of the metaplectic algebra}
    \label{ssec:lorentzian-complexification}

    The metaplectic lift is real and symplectic.
    Its Lie algebra is $\mathfrak{sp}(2,\mathbb{R}) = \mathfrak{sl}(2,\mathbb{R})$, since the double cover
    does not change the Lie algebra, and its complexification is
    \begin{equation}
      \label{eq:metaplectic-complexification}
      \mathfrak{mp}(2,\mathbb{R})_{\mathbb{C}} \simeq \mathfrak{sl}_2(\mathbb{C}).
    \end{equation}
    This is a statement about Lie algebras and uses no metric.
    The simply connected Lie group integrating $\mathfrak{sl}_2(\mathbb{C})$ is $\SL(2,\mathbb{C})$, and as an
    abstract Lie group
    \begin{equation}
      \label{eq:spin-group-lorentzian}
      \Spin(3,1) \simeq \SL(2,\mathbb{C}).
    \end{equation}
    An abstract isomorphism with $\Spin(3,1)$ does not make this group the spin group of the effective
    spacetime.
    That requires a Lorentzian co-metric on the effective base and an identification of the
    $\SL(2,\mathbb{C})$ acting on the admissible doublet with the group covering the Lorentz frames of that
    co-metric: Hypothesis~\ref{hyp:spin-solder} below.

    Let $V \simeq \mathbb{C}^2$ denote the admissible doublet, the two-dimensional carrier on which the
    Born--Infeld datum of Subsection~\ref{ssec:proof-bi-left} acts; it is supplied model data, not derived
    from the Weil module~\cite{Beau2026a27,Beau2026a34}.
    It carries the fundamental representation of $\SL(2,\mathbb{C})$ through~\eqref{eq:metaplectic-complexification}.

  \subsection{The admissible spin frame bundle under [H-Spin]}
    \label{ssec:spin-frame-bundle}

    Under Hypothesis~\ref{hyp:spin-solder}, let $P_\Pi$ denote the principal $\Spin(3,1)$-bundle it
    supplies.
    The left-handed and right-handed Weyl bundles are
    \[
      S_L = P_\Pi \times_{\Spin(3,1)} V,
      \qquad
      S_R = P_\Pi \times_{\Spin(3,1)} \overline{V},
    \]
    and the admissible spinor bundle is the Dirac bundle $S_\Pi = S_L \oplus S_R$.
    The construction is projective in intent: $P_\Pi$ is meant to be induced from the Weil fibre rather
    than postulated as a spin frame bundle over a pre-existing manifold.
    Hypothesis~\ref{hyp:spin-solder} is exactly the statement that this induction exists, and no source
    supplies it.

  \subsection{Tensor sectors: the Lorentz reading and the electroweak reading}
    \label{ssec:electroweak-tensors}

    On the abstract doublet the tensor identities are
    \begin{equation}
      \label{eq:sym2-fibre}
      \operatorname{Sym}^2(V) \simeq \mathfrak{sl}_2(\mathbb{C}),
      \qquad
      \wedge^2(V) \simeq \mathbb{C}_{\det} ,
    \end{equation}
    where $\operatorname{Sym}^2(V)$ carries the adjoint representation and $\wedge^2(V)$ the trivial one,
    since $\det g = 1$ on $\SL(2,\mathbb{C})$.
    A Hermitian form on $V$ selects a compact real form $\mathfrak{su}(2) \subset \mathfrak{sl}_2(\mathbb{C})$.

    These identities admit two readings, which must not be taken for one copy of $V$.
    In the Lorentz reading, under Hypothesis~\ref{hyp:spin-solder}, $V$ is the Weyl fibre of $S_L$ and
    \begin{equation}
      \label{eq:sym2-adjoint}
      \operatorname{Sym}^2(S_L) \simeq \operatorname{ad}(P_\Pi) ,
    \end{equation}
    the bundle of Lorentz Lie algebras with its complex structure, that is the complex self-dual two-forms.
    It is a Lorentz sector, not a weak adjoint.
    The line $\wedge^2(S_L)$ is trivial, spanned by the invariant form $\varepsilon$ of
    Lemma~\ref{lem:covariant-epsilon}, and supplies no non-trivial hypercharge.
    A compact real form of $\mathfrak{sl}_2(\mathbb{C})$ is then the stabiliser of a time-like direction: a
    positive Hermitian form on $S_L$ is a time-like vector, and the corresponding $\mathfrak{su}(2)$
    generates spatial rotations in that frame, not an internal symmetry.

    In the electroweak reading the same identities are applied to a second rank-two factor.
    On an irreducible Weyl fibre, every complex-linear map commuting with the Lorentz action is scalar
    (Schur's lemma), so a non-trivial weak $\SU(2)$ action commuting with Lorentz transformations requires
    an additional space: the field carries two indices,
    \[
      \psi_{\alpha i} \in S_L \otimes E_{\mathrm{weak}},
      \qquad
      \rho(\Lambda, u) = \rho_L(\Lambda) \otimes \rho_{\mathrm{weak}}(u),
    \]
    as in the Standard Model~\cite{Tong2024SM}, and as diagnosed for this programme in the spin-rotation
    note~\cite{Beau2026srn}.
    The factor $E_{\mathrm{weak}}$ is Hypothesis~\ref{hyp:weak-factor}.
    Under it, $\operatorname{Sym}^2$ applied to the fibre of $E_{\mathrm{weak}}$ gives the complex weak
    adjoint, whose compact real form selected by the Hermitian structure of $E_{\mathrm{weak}}$ is the
    $\SU(2)_L$ algebra, and $L_Y := \wedge^2(E_{\mathrm{weak}})$ is the determinant line.
    The line $L_Y$ carries a non-trivial $U(1)_Y$ action only because the structure group of
    $E_{\mathrm{weak}}$ is taken to be $U(2)$ rather than $\SU(2)$; that choice is part of
    Hypothesis~\ref{hyp:weak-factor}.
    It fixes the existence of the hypercharge line, while the distribution of hypercharge weights is
    constrained by spectral anomaly cancellation in Section~\ref{sec:anomaly}.

    \begin{theorem}[Spinorial electroweak bundle]
      \label{thm:spinorial-electroweak-bundle}
      \begin{enumerate}
        \item[\emph{(a)}] The complexified metaplectic algebra of the admissible Weil module
        satisfies~\eqref{eq:metaplectic-complexification}, integrates to $\SL(2,\mathbb{C})$, and on the
        abstract doublet $V$ the tensor identities~\eqref{eq:sym2-fibre} hold; a Hermitian form on $V$
        selects the compact real form $\mathfrak{su}(2)$.
        \item[\emph{(b)}] Under Hypothesis~\ref{hyp:spin-solder}, $S_\Pi = S_L \oplus S_R$ is a Dirac spinor
        bundle of the effective Lorentzian co-metric, $\operatorname{Sym}^2(S_L) \simeq
        \operatorname{ad}(P_\Pi)$ is its Lorentz sector, and $\wedge^2(S_L)$ is trivial.
        \item[\emph{(c)}] Under Hypothesis~\ref{hyp:weak-factor}, the identities~\eqref{eq:sym2-fibre}
        applied to $E_{\mathrm{weak}}$ give the complex weak adjoint, with compact real form
        $\mathfrak{su}(2)_L$, and the determinant line $L_Y = \wedge^2(E_{\mathrm{weak}})$ carrying
        $U(1)_Y$; together these are the $\SU(2)_L \times U(1)_Y$ electroweak bundle data on
        $S_\Pi \otimes E_{\mathrm{weak}}$.
      \end{enumerate}
      Part~(a) is proved without either hypothesis.
    \end{theorem}

    \begin{proof}
      (a) The metaplectic group is a double cover of $\mathrm{Sp}(2,\mathbb{R}) = \SL(2,\mathbb{R})$, so its Lie
      algebra is $\mathfrak{sl}(2,\mathbb{R})$, whose complexification is $\mathfrak{sl}_2(\mathbb{C})$; the
      simply connected group with this Lie algebra is $\SL(2,\mathbb{C})$.
      For the fundamental representation, $\operatorname{Sym}^2(\mathbb{C}^2)$ is the three-dimensional
      irreducible representation, which is the adjoint, and $\wedge^2(\mathbb{C}^2)$ is the determinant,
      trivial on $\SL(2,\mathbb{C})$.
      The unitary group of a positive Hermitian form on $\mathbb{C}^2$ meets $\SL(2,\mathbb{C})$ in a conjugate
      of $\SU(2)$, whose Lie algebra is a compact real form of $\mathfrak{sl}_2(\mathbb{C})$.
      (b) Under Hypothesis~\ref{hyp:spin-solder}, the associated-bundle construction applied to the fibrewise
      identities gives $\operatorname{Sym}^2(S_L) \simeq P_\Pi \times_{\mathrm{Ad}} \mathfrak{sl}_2(\mathbb{C})
      = \operatorname{ad}(P_\Pi)$ and $\wedge^2(S_L) \simeq P_\Pi \times_{\det} \mathbb{C}$, which is trivial.
      (c) Under Hypothesis~\ref{hyp:weak-factor}, the same construction applied to the $U(2)$-bundle
      $E_{\mathrm{weak}}$ gives $\operatorname{Sym}^2(E_{\mathrm{weak}})$, containing the complexified
      $\mathfrak{su}(2)_L$ adjoint, and $\wedge^2(E_{\mathrm{weak}}) = \det E_{\mathrm{weak}}$, on which
      $U(2)$ acts through the determinant character.
    \end{proof}

    Fermions are read from the Weil module through the complexification of its metaplectic Lie algebra.
    The algebra is proved; its reading as spacetime spinors and its reading as weak doublets are two
    distinct conditions, stated next.

    \begin{hypothesis}[{[H-Spin]}: Lorentzian spin solder; not supplied]
      \label{hyp:spin-solder}
      The following two conditions hold.
      \begin{enumerate}
        \item[\emph{(i)}] The effective base $M_{\mathrm{eff}}$ carries a four-dimensional Lorentzian co-metric
        $g$, time- and space-orientable.
        \item[\emph{(ii)}] There is a principal $\Spin(3,1)$-bundle $P_\Pi \to M_{\mathrm{eff}}$ with a
        two-to-one equivariant map onto the bundle of oriented, time-oriented orthonormal frames of $g$,
        such that the $\SL(2,\mathbb{C})$ acting on the admissible doublet $V$
        through~\eqref{eq:metaplectic-complexification} is the structure group of $P_\Pi$.
      \end{enumerate}
    \end{hypothesis}

    Part~(i) is supplied only conditionally by the geometric branch.
    Under the spatial limit hypothesis [H-L] and the lifting hypothesis [H-lift], the effective co-metric
    of Q5b has rank three, of signature $(-,+,+)$ under the reduced hyperbolicity hypothesis
    [H-hyp$_{\mathrm{red}}$], and a full-rank extension is open (Q5b~\cite{Beau2026q5b}, Theorems~5.2
    and~6.1, open problem Q5b-O2).
    For a supplied full-rank extension with a spatially $\SU(2)$-invariant principal symbol and a negative
    ordering term, Q11 obtains $g^{\mu\nu} = \mathrm{diag}(-A_\tau, 2\lambda, 2\lambda, 2\lambda)$ with
    $A_\tau$ and $\lambda$ independent (Q11~\cite{Beau2026q11}, Corollary~6.1, using
    Q8~\cite{BeauQ8}); no coefficient value is derived, and only the signature is used here.
    Part~(ii) is supplied by no source: the metaplectic action is an action on the internal Weil fibre,
    and no result of the corpus relates it to the frames of $g$.
    Neither part is excluded, and Hypothesis~\ref{hyp:spin-solder} is a missing identification: the
    Lorentz-covariant statements of this paper hold as stated once it is supplied.

    \begin{hypothesis}[{[H-Weak]}: distinct weak factor; not supplied]
      \label{hyp:weak-factor}
      There is a Hermitian vector bundle $E_{\mathrm{weak}} \to M_{\mathrm{eff}}$ of rank two, with structure
      group $U(2)$ and a connection, distinct from the factor soldered by
      Hypothesis~\ref{hyp:spin-solder}, and the fermion fields are sections of the twisted bundle
      $S_\Pi \otimes E_{\mathrm{weak}}$ with the representation content of the Standard Model: left-handed
      components in doublets of $E_{\mathrm{weak}}$ and right-handed components in weak singlets.
    \end{hypothesis}

    This paper does not construct $E_{\mathrm{weak}}$.
    By Schur's lemma, the admissible doublet can supply at most one of the two rank-two factors: the chain
    of~\eqref{eq:main-chain} assigns it to the spinor factor, through Hypothesis~\ref{hyp:spin-solder}, and
    reading it instead as the weak factor would require the spinor bundle to come from elsewhere.
    Which reading, if either, the admissible Weil fibre realises is not decided here.

  \section{The Projected Dirac Operator and the Projective Endomorphism}
  \label{sec:dirac}

  The preceding section constructs the admissible spinor bundle $S_\Pi$ under
  Hypothesis~\ref{hyp:spin-solder}.
  We now define the Dirac operator carried by this bundle, isolate the endomorphism
  induced by non-injective projection, and establish its role in chiral selection.
  Throughout Sections~\ref{sec:dirac} and~\ref{sec:anomaly}, Hypothesis~\ref{hyp:spin-solder} is
  assumed, and the gauge connection acts on the weak factor of Hypothesis~\ref{hyp:weak-factor}.
  When the connection is included, the Dirac operator acts on the twisted bundle
  $S_\Pi \otimes E_{\mathrm{weak}}$ (with the appropriate hypercharge and colour factors); we keep the
  notation $S_\Pi$ for it when no confusion arises, the Clifford matrices and $\gamma_5$ acting on the
  spinor factor and the gauge generators $T^a_R$ and the hypercharge $Y$ on the internal factors.

  \subsection{The minimal spin--gauge Dirac operator}
    \label{ssec:minimal-dirac}

    Let $\nabla^{S,A}$ denote the spin--gauge connection on $S_\Pi$, obtained by
    combining the Levi--Civita spin connection of the effective metric $g$ with the
    admissible gauge connection $A$ of~\cite{Beau2026q6a,Beau2026q12}.
    The minimally coupled Dirac operator is
    \begin{equation}
      \label{eq:minimal-dirac}
      \mathcal{D}_{g,A} = \gamma^\mu \nabla^{S,A}_\mu .
    \end{equation}
    In an ordinary $\mathrm{spin}^c$ geometry over a given manifold, the square of
    this operator is governed by the standard Lichnerowicz formula.
    In the present setting, $S_\Pi$ is not a primitive spinor bundle over a pre-existing
    manifold: it is the spinorial face of the admissible Weil fibre, and the projection
    $\Pi$ acts non-trivially on the spinorial sector.

  \subsection{The projective correction and the definition of
  \texorpdfstring{$E_\Pi$}{E\_Pi}}
    \label{ssec:projective-correction}

    Let $\Pi_S \colon S \to S_\Pi$ denote the restriction of the admissible projection
    to the spinorial bundle.
    The projected Dirac operator is defined by
    \begin{equation}
      \label{eq:projected-dirac}
      \mathcal{D}_{\Pi,g,A} := \Pi_S \circ \mathcal{D}_{g,A} \circ \Pi_S^* ,
    \end{equation}
    where $\Pi_S^*$ denotes the admissible minimal lift, equivalently the adjoint of
    $\Pi_S$ after choosing the induced Hermitian structure on the projected spinor
    bundle.
    We do not assume that $\Pi_S \Pi_S^* = \mathrm{id}_{S_\Pi}$; this identity holds
    precisely in the injectively resolved limit and is in general replaced by a
    non-trivial projector.
    We assume throughout this section that the admissible lift is symbol-compatible:
    the principal symbol of $\mathcal{D}_{\Pi,g,A}$ coincides with the projected
    Clifford symbol $\gamma^\mu \xi_\mu$.
    Equivalently, the projection defect affects only the zero-order part of
    $\mathcal{D}_{\Pi,g,A}^2$, not its principal symbol.

    \begin{definition}[Projective endomorphism]
      \label{def:projective-endomorphism}
      The projective endomorphism $E_\Pi \in \Gamma(\operatorname{End}(S_\Pi))$ is
      defined by
      \begin{equation}
        \label{eq:epi-definition}
        E_\Pi
        :=
        \mathcal{D}_{\Pi,g,A}^2
        -
        \left(
          -(\nabla^{S,A})^2
        + \frac{R}{4}
          + \frac{1}{2}\gamma^\mu\gamma^\nu F^a_{\mu\nu} T^a_R
        \right).
      \end{equation}
    \end{definition}

    The right-hand side of~\eqref{eq:epi-definition} collects the terms of the
    standard Lichnerowicz formula.
    The endomorphism $E_\Pi$ is therefore not an additional matter potential.
    It is the internal spectral residue of non-injective projection in the spinorial
    sector: the part of $\mathcal{D}_{\Pi,g,A}^2$ not resolved by the effective
    pair $(g, A)$.

    \begin{lemma}[Projective origin of $E_\Pi$]
      \label{lem:projective-endomorphism}
      Let $S_\Pi$ be the admissible spinor bundle of Hypothesis~\ref{hyp:spin-solder}, and assume
      that the admissible lift $\Pi_S^*$ is symbol-compatible with the spin--gauge
      connection.
      The projected square $\mathcal{D}_{\Pi,g,A}^2$ then differs from the ordinary
      spin--gauge Lichnerowicz square by a canonical zero-order endomorphism
      $E_\Pi \in \Gamma(\operatorname{End}(S_\Pi))$.
      If the spinorial projection is injectively resolved by the effective pair $(g,A)$,
      then $E_\Pi = 0$.
      Conversely, $E_\Pi \neq 0$ whenever the non-injective spinorial fibres leave a
      non-trivial zero-order trace in the projected square.
      Thus axiom~A2 of~\cite{Beau2026m} makes $E_\Pi$ the natural place where
      structural non-injectivity enters the spinorial spectral action.
    \end{lemma}

    \begin{proof}
      When $\Pi_S$ is injective, $\Pi_S^*$ is its right inverse and
      $\Pi_S \Pi_S^* = \mathrm{id}_{S_\Pi}$.
      The composition~\eqref{eq:projected-dirac} then reduces to $\mathcal{D}_{g,A}$
      on $S_\Pi$, whose square is the standard Lichnerowicz formula; hence $E_\Pi = 0$.
      When $\Pi_S \Pi_S^* \neq \mathrm{id}_{S_\Pi}$, the squared projected operator
      contains additional terms controlled by the projection defect.
      By symbol-compatibility, these terms do not modify the principal Clifford symbol
      and therefore contribute only to the zero-order part of the Laplace-type operator.
      They are collected canonically in $E_\Pi$ via
      Definition~\ref{def:projective-endomorphism}.
    \end{proof}

    \begin{remark}
      \label{rem:symbol-compatibility-seeley-dewitt}
      Symbol-compatibility is also required in Section~\ref{sec:anomaly}: the computation
      of the Seeley--DeWitt coefficients of $\mathcal{D}_{\Pi,g,A}^2$ via the standard
      heat-kernel expansion presupposes that the principal symbol is the Clifford symbol,
      so that the leading heat-kernel term is the standard Gaussian.
      The hypothesis is therefore not an isolated technical convenience, but a coherence
      condition between the projective construction and the spectral action formalism.
    \end{remark}

  \subsection{Comparison with the Connes--Chamseddine approach}
    \label{ssec:comparison-cc}

    This is the point at which the present construction differs structurally from the
    Connes--Chamseddine spectral action~\cite{ConnesChamseddine1997}.
    In the almost-commutative setting, the finite internal Dirac operator is part of
    the input data of the spectral triple, and Yukawa matrices and fermion masses are
    encoded directly in this finite component.
    Here $E_\Pi$ is not postulated as a finite Dirac datum: it is induced by the
    non-injective projection itself.
    The task is therefore not to choose $E_\Pi$, but to determine which endomorphisms
    are compatible with admissibility, gauge invariance, and spectral coherence.
    This inversion of logical order is the distinctive feature of the present approach:
    internal fermionic structure is not input but output.

  \subsection{Chiral admissibility and the \texorpdfstring{$V{-}A$}{V-A} structure}
    \label{ssec:chiral-admissibility}

    Let $\gamma_5 = i \gamma^0 \gamma^1 \gamma^2 \gamma^3$ be the chirality operator
    of $\mathrm{Cl}(T^* M_{\mathrm{eff}}, g^{\mu\nu})$, whose existence as a global
    section depends on the oriented Lorentzian spin structure of Hypothesis~\ref{hyp:spin-solder}.
    Define the chiral projectors
    \[
      P_L = \frac{1}{2}(1 - \gamma_5),
      \qquad
      P_R = \frac{1}{2}(1 + \gamma_5).
    \]

    \begin{definition}[Chiral projective admissibility]
      \label{def:chiral-admissibility}
      The projective endomorphism $E_\Pi$ is called \emph{left-admissible} if
      \begin{equation}
        \label{eq:left-admissible}
        P_R \, E_\Pi \, P_R = 0 .
      \end{equation}
      Equivalently, the purely right-handed self-coupling of the projective residue
      vanishes.
    \end{definition}

  \subsection{Spinorial lift of BI parity}
    \label{ssec:proof-bi-left}

    The spinorial lift is derived here from a supplied Born--Infeld datum.
    The Weil-action conjugation identity $\rho_{q-c} = \bar\rho_c$~\cite{Beau2026a22} relates the
    sectors $c$ and $q-c$; turning it into an antilinear involution of one admissible doublet requires
    an inter-block identification and a weight exchange, which are supplied as hypotheses (H1) and (H2)
    of O30~\cite{Beau2026a34} and are not derived there.
    Under these hypotheses the parity is represented by an antilinear involution
    \[
      K : V_\rho \longrightarrow V_\rho,
      \qquad
      K^2 = 1 ,
    \]
    on the doublet $V_\rho = V_+ \oplus V_-$ of Section~\ref{sec:metaplectic}.
    No convergence estimate of the finite-$q$ Weil blocks toward $K$ is supplied by the corpus.
    By (H2) the involution exchanges the two weight sectors, and the grading
    \[
      H := 2 J_3,
      \qquad
      H|_{V_+} = +1,
      \qquad
      H|_{V_-} = -1,
      \qquad
      H^2 = 1 ,
    \]
    records which sector a vector belongs to.

    \begin{lemma}[BI parity is off-diagonal with respect to its grading]
      \label{lem:bi-grading-anticommute}
      The limiting BI involution $K$ and the weight grading $H$ satisfy
      \begin{equation}
        \label{eq:HK-anticommute}
        H K = - K H .
      \end{equation}
    \end{lemma}

    \begin{proof}
      For $v_+ \in V_+$ one has $K v_+ \in V_-$, hence $H K v_+ = - K v_+$ and $K H v_+ = K v_+$,
      so $H K v_+ = - K H v_+$.
      The same computation on $V_-$ gives the operator identity.
    \end{proof}

    Equation~\eqref{eq:HK-anticommute} is the operator expression of the fact that the BI parity
    exchanges precisely the two weights that the supplied BI datum grades.
    Lemma~\ref{lem:bi-grading-anticommute} and its consequence $(HK)^2 = -1$ use no metric; the rest
    of this subsection assumes Hypothesis~\ref{hyp:spin-solder}.

    \begin{lemma}[Unique covariant antilinear structure]
      \label{lem:covariant-epsilon}
      Assume Hypothesis~\ref{hyp:spin-solder}.
      Then the admissible doublet is the Weyl fibre $S_L \simeq \mathbb{C}^2$ of the Lorentz spin group
      $\SL(2,\mathbb{C})$, with $S_R \simeq \overline{S_L}$ inequivalent to $S_L$.
      Consequently no Lorentz-covariant antilinear conjugation $S_L \to S_L$ exists, so any
      covariant antilinear lift must exchange the chiral subbundles
      $P_L S_\Pi \leftrightarrow P_R S_\Pi$, and the unique $\SL(2,\mathbb{C})$-invariant bilinear
      form on $S_L$ is, up to scale,
      \[
        \varepsilon =
        \begin{pmatrix} 0 & 1 \\ -1 & 0 \end{pmatrix},
        \qquad
        \varepsilon^2 = -1 .
      \]
    \end{lemma}

    \begin{proof}
      A covariant antilinear endomorphism of $S_L$ is a covariant isomorphism
      $S_L \simeq \overline{S_L}$, that is $\rho \simeq \overline{\rho}$ for the fundamental
      representation; the $\SL(2,\mathbb{C})$ Weyl representations $(\tfrac12,0)$ and $(0,\tfrac12)$
      are inequivalent, so no such map exists, which forces the off-diagonal exchange.
      For the form, invariance $g^{T} B g = B$ under the diagonal torus forces the diagonal entries
      of $B$ to vanish, and invariance under the unipotent forces the off-diagonal entries to be
      opposite, so $B \propto \varepsilon$, which is itself invariant under $\SL(2,\mathbb{C})$.
      Under the compact form $\mathfrak{su}(2)$ the doublet is pseudoreal~\cite{FultonHarris1991}; it is
      the reading of $\SL(2,\mathbb{C})$ as the Lorentz spin group, Hypothesis~\ref{hyp:spin-solder},
      that removes the equivalence $\rho \simeq \overline{\rho}$ and converts the internal weight exchange
      into a chiral exchange.
    \end{proof}

    \begin{theorem}[Spinorial BI lift]
      \label{thm:spinorial-bi-lift}
      Assume Hypothesis~\ref{hyp:spin-solder} and the Born--Infeld datum $(K,H)$ supplied by
      (H1)--(H2) of~\cite{Beau2026a34}.
      On the spinor bundle $S_\Pi$ there exists a unique antilinear lift of the BI parity
      compatible with both the BI datum $(K,H)$ and $\SL(2,\mathbb{C})$-covariance.
      Its internal part is the grading-dressed parity $J_\Pi^{\mathrm{int}} = H K$, which in the
      weight basis ($K = \sigma_x \circ \overline{(\cdot)}$, $H = \sigma_z$) equals the unique
      covariant form $\varepsilon \circ \overline{(\cdot)}$, since $\sigma_z \sigma_x = \varepsilon$.
      The lift satisfies
      \begin{equation}
        \label{eq:jpi-conditions}
        J_\Pi^2 = -1,
        \qquad
        J_\Pi \gamma_5 = -\gamma_5 J_\Pi .
      \end{equation}
    \end{theorem}

    \begin{proof}
      By Lemma~\ref{lem:covariant-epsilon} the lift is off-diagonal with internal part proportional
      to $\varepsilon \circ \overline{(\cdot)}$, and compatibility with the BI datum fixes this
      internal part to $H K$.
      This is Lorentz-legitimate although $H$ alone is not invariant, precisely because $H K$
      coincides with the unique invariant $\varepsilon \circ \overline{(\cdot)}$ of
      Lemma~\ref{lem:covariant-epsilon}; the same two constraints determine the lift uniquely.
      The quaternionic sign is then forced rather than imposed:
      \[
        (H K)^2 = H K H K = - H^2 K^2 = -1 ,
      \]
      using~\eqref{eq:HK-anticommute}, $H^2 = 1$, and $K^2 = 1$.
      The first identity in~\eqref{eq:jpi-conditions} follows from $(H K)^2 = -1$; the second from
      the off-diagonal exchange of the chiral subbundles, which reverses chirality while remaining
      involutive and therefore contributes no sign to the square.
    \end{proof}

    The first condition in~\eqref{eq:jpi-conditions} expresses the quaternionic
    character of the spinorial lift.
    The second states that $J_\Pi$ reverses chirality.
    Together they imply
    \begin{equation}
      \label{eq:jpi-chiral-exchange}
      J_\Pi P_L = P_R J_\Pi,
      \qquad
      J_\Pi P_R = P_L J_\Pi,
    \end{equation}
    so the BI parity lift exchanges the two chiral subbundles $P_L S_\Pi$ and
    $P_R S_\Pi$.

    \begin{proposition}[BI parity and left-admissibility]
      \label{prop:bi-left-admissibility}
      By Theorem~\ref{thm:spinorial-bi-lift}, and assuming that admissibility selects the
      orientation-compatible branch, $E_\Pi$ is left-admissible:
      \[
        P_R E_\Pi P_R = 0 .
      \]
    \end{proposition}

    \begin{proof}
      By~\eqref{eq:jpi-chiral-exchange}, any right-handed self-coupling of a section
      of $S_\Pi$ is mapped under $J_\Pi$ to a left-handed coupling.
      The admissibility condition selects the branch compatible with the orientation
      carried by $J_\Pi$; equivalently, projective self-couplings of $P_R S_\Pi$ are
      not orientation-compatible with the BI parity involution.
      Hence $P_R E_\Pi P_R = 0$.
    \end{proof}

    On the left-admissible branch, the projective residue is chiral block-diagonal.
    The Schur form $E_\Pi = -M^\dagger M \preceq 0$ established in~\cite{Beau2026prs},
    together with left-admissibility $P_R E_\Pi P_R = 0$, forces the off-diagonal
    chiral blocks to vanish: a self-adjoint negative-semidefinite operator with a
    vanishing diagonal block has vanishing off-diagonal blocks.
    Consequently
    \begin{equation}
      \label{eq:epi-gamma5-commutator}
      [E_\Pi, \gamma_5] = 0 ,
    \end{equation}
    and $E_\Pi = P_L E_\Pi P_L$.
    The chiral asymmetry of the projective residue is therefore carried not by a chiral
    commutator but by the diagonal imbalance between $P_L E_\Pi P_L$ and
    $P_R E_\Pi P_R = 0$, equivalently by the $\gamma_5$-graded trace, which enters the spectral
    anomaly condition of Section~\ref{sec:anomaly}.

    \begin{remark}[Lorentz chirality and electroweak chiral selection]
      \label{rem:lorentz-vs-weak-chirality}
      Proposition~\ref{prop:bi-left-admissibility} is a statement about the projective residue $E_\Pi$
      and the Lorentz chirality $\gamma_5$ of the spinor factor; it is conditional on
      Hypothesis~\ref{hyp:spin-solder}, the supplied Born--Infeld datum and the orientation branch.
      The chiral selection of the weak interaction is a different statement: that the weak gauge action
      is carried by the left-handed factor alone, with fermions in
      $P_L S_\Pi \otimes E_{\mathrm{weak}}$ and right-handed components in weak singlets.
      That is the representation content of Hypothesis~\ref{hyp:weak-factor}.
      No argument in this paper derives it from $P_R E_\Pi P_R = 0$; relating the two would require a
      coupling between the projective residue and the weak factor, which is not constructed.
      The $V{-}A$ structure of weak couplings therefore rests on Hypothesis~\ref{hyp:weak-factor} with
      that chiral assignment, and left-admissibility is compatible with it but does not select it.
    \end{remark}

  \section{Spectral Anomaly Cancellation and Hypercharge Rigidity}
  \label{sec:anomaly}

  The preceding section shows that $\mathcal{D}_{\Pi,g,A}^2$ is a Laplace-type
  operator with standard Clifford principal symbol and projective endomorphism $E_\Pi$.
  We now use this structure to formulate the hypercharge consistency condition.
  The guiding principle is that the projected spectral functional
  $S_\Pi[g, A, \mathcal{D}_{\Pi,g,A}]$ must be invariant under $U(1)_Y$ gauge
  transformations.
  In the spinorial sector, this invariance imposes trace constraints on the
  hypercharge weights carried by the projected components of $E_\Pi$.
  These constraints are the spectral form of anomaly cancellation.

  \subsection{Projected \texorpdfstring{$U(1)_Y$}{U(1)\_Y} variation}
    \label{ssec:u1-variation}

    Let $L_Y = \wedge^2(E_{\mathrm{weak}})$ be the determinant line identified in
    Section~\ref{sec:metaplectic}.
    The projected spinorial sector decomposes into irreducible components
    $S_R \subset S_\Pi$, each carrying a definite chirality and hypercharge weight
    $Y_R$.
    A $U(1)_Y$ gauge transformation acts by
    \begin{equation}
      \label{eq:u1-action}
      \psi_R \longmapsto e^{i Y_R \alpha} \psi_R ,
    \end{equation}
    and the infinitesimal action is generated by the endomorphism
    $Y = \bigoplus_R Y_R \, \mathrm{id}_{S_R}$.
    The projected spectral functional is hypercharge-consistent if
    \begin{equation}
      \label{eq:u1-invariance}
      \delta_{\alpha Y} S_\Pi[g, A, \mathcal{D}_{\Pi,g,A}] = 0
    \end{equation}
    for all smooth gauge parameters $\alpha$.

  \subsection{The \texorpdfstring{$\gamma_5$-weighted $a_4$}{gamma5-weighted a4}
  coefficient}
    \label{ssec:gamma5-a4}

    By symbol-compatibility (Section~\ref{ssec:projective-correction}),
    $\mathcal{D}_{\Pi,g,A}^2$ is a Laplace-type operator with Clifford principal symbol,
    written in standard form as
    \begin{equation}
      \label{eq:laplace-form}
      \mathcal{D}_{\Pi,g,A}^2 = -(\nabla^{S,A})^2 + E,
    \end{equation}
    with
    \begin{equation}
      \label{eq:endomorphism-E}
      E = \frac{R}{4} + \frac{1}{2}\gamma^\mu\gamma^\nu F^a_{\mu\nu} T^a_R + E_\Pi .
    \end{equation}
    In four dimensions, the local chiral anomaly is controlled by the
    $\gamma_5$-weighted $a_4$ coefficient of the Seeley--DeWitt expansion.
    We denote the corresponding local endomorphism-valued anomaly density by
    $\mathcal{A}_\Pi$.
    After taking the spinorial trace, it defines the scalar anomaly density entering
    the variation of the projected spectral functional.
    It is determined by the curvature of $\nabla^{S,A}$, the projective endomorphism
    $E_\Pi$, and the effective metric $g$.

    The local chiral anomaly density is
    \begin{equation}
      \label{eq:anomaly-density-a4}
      \mathcal{A}_\Pi(x) = a_4\!\left(x,\, \mathcal{D}_{\Pi,g,A}^2;\, \gamma_5\right).
    \end{equation}

    \begin{lemma}[Spectral trace form of the anomaly]
      \label{lem:spectral-anomaly-trace}
      Assume symbol-compatibility of the projected Dirac operator.
      Then the anomalous $U(1)_Y$ variation of the projected spectral functional is
      \begin{equation}
        \label{eq:anomaly-trace}
        \delta_{\alpha Y} S_\Pi
        =
        \int_{M_{\mathrm{eff}}}
        \alpha \,
        \operatorname{Tr}_{S_\Pi}\!\left(
                                     \gamma_5 \, Y \, \mathcal{A}_\Pi(x)
        \right)
        \sqrt{|g|} \, d^4x ,
      \end{equation}
      where $\mathcal{A}_\Pi$ is the local endomorphism-valued $a_4$ Seeley--DeWitt
      density of $\mathcal{D}_{\Pi,g,A}^2$.
      Gauge invariance therefore requires the local vanishing condition
      \begin{equation}
        \label{eq:anomaly-vanishing}
        \operatorname{Tr}_{S_\Pi}\!\left(
                                     \gamma_5 \, Y \, \mathcal{A}_\Pi(x)
        \right) = 0
        \qquad
        \text{for all } x \in M_{\mathrm{eff}} .
      \end{equation}
    \end{lemma}

    \begin{remark}
      \label{rem:seeley-dewitt-hierarchy}
      The appearance of the $\gamma_5$-weighted $a_4$ coefficient is consistent with
      the spectral stratification proposed, as an interpretive reading, in the bosonic
      sector~\cite{Beau2026q13}.
      In the bosonic variation, the $a_2$ coefficient gives the gravitational response
      and the $a_4$ coefficient gives the gauge response.
      In the spinorial sector, the same $a_4$ level controls the local chiral anomaly
      and therefore the hypercharge consistency of the fermionic gauge coupling.
      Anomaly cancellation is therefore not a new spectral layer: it is the chiral
      consistency condition of the gauge-level response, located at exactly the same
      Seeley--DeWitt order as Yang--Mills dynamics.
      Schematically:
      \begin{equation}
        \label{eq:a4-two-faces}
        a_4 \;\longrightarrow\;
        \begin{cases}
          \delta_A S_\Pi \Rightarrow \text{Yang--Mills dynamics}, \\[4pt]
          \delta_g S_\Pi\big|_{a_4} \Rightarrow \text{Yang--Mills stress tensor}, \\[4pt]
          \operatorname{Tr}_{S_\Pi}(\gamma_5 \, \delta_{\alpha Y}) \Rightarrow
          \text{anomaly cancellation}.
        \end{cases}
      \end{equation}
    \end{remark}

  \subsection{Trace constraints on projected hypercharge weights}
    \label{ssec:trace-constraints}

    Only terms in $a_4$ containing four Clifford matrices survive the $\gamma_5$-trace.
    The gauge-curvature part of~\eqref{eq:endomorphism-E} contributes the topological
    density
    \begin{equation}
      \label{eq:topological-density}
      \operatorname{Tr}_{S_\Pi}\!\left(
                                   \gamma_5 Y \, \gamma^\mu\gamma^\nu\gamma^\rho\gamma^\sigma
                                   F_{\mu\nu} F_{\rho\sigma}
      \right)
      \propto \operatorname{tr}_R(Y F \wedge F).
    \end{equation}
    For the abelian factor this gives the cubic $U(1)_Y$ anomaly; for non-abelian
    factors it gives the mixed $U(1)_Y \times G^2$ anomalies.
    The projective endomorphism $E_\Pi$ contributes additional chiral traces through
    the chiral imbalance of its diagonal blocks, with $E_{RR} = 0$ by left-admissibility,
    as read by the $\gamma_5$-graded trace.
    Left-admissibility, $P_R E_\Pi P_R = 0$, constrains these contributions to the
    orientation-compatible sector.
    They are therefore absorbed into the same projected anomaly condition, without
    generating independent constraints beyond those arising from the gauge-curvature
    sector.

    Let $\chi_R \in \{+1, -1\}$ denote the chirality eigenvalue of sector $R$:
    $\chi_R = +1$ for left-handed sectors ($S_R \subset P_L S_\Pi$) and
    $\chi_R = -1$ for right-handed sectors ($S_R \subset P_R S_\Pi$), as fixed by the chiral
    assignment of Hypothesis~\ref{hyp:weak-factor} (Remark~\ref{rem:lorentz-vs-weak-chirality}).
    The vanishing condition~\eqref{eq:anomaly-vanishing} specialises, at leading
    Seeley--DeWitt order, to the following set of spectral coherence constraints:
    \begin{align}
      \label{eq:anomaly-linear}
      \sum_R \chi_R \, Y_R \, \dim(R) &= 0 , \\
      \label{eq:anomaly-cubic}
      \sum_R \chi_R \, Y_R^3 \, \dim(R) &= 0 , \\
      \label{eq:anomaly-mixed}
      \sum_R \chi_R \, Y_R \, \operatorname{tr}_R(T^a T^b) &= 0 .
    \end{align}
    Equations~\eqref{eq:anomaly-linear}--\eqref{eq:anomaly-mixed} are the
    gravitational, cubic $U(1)_Y$, and mixed $U(1)_Y \times G^2$ anomaly conditions,
    respectively.
    They are not imposed externally; they follow from requiring that the projected
    spectral functional be well-defined under $U(1)_Y$.

    \begin{proposition}[Spectral anomaly constraints]
      \label{prop:spectral-anomaly-constraints}
      Assume Hypotheses~\ref{hyp:spin-solder} and~\ref{hyp:weak-factor}, symbol-compatibility, and
      left-admissibility of $E_\Pi$ (Theorem~\ref{thm:spinorial-bi-lift},
      Proposition~\ref{prop:bi-left-admissibility}).
      Then $U(1)_Y$ invariance of the projected spectral functional imposes the
      projected anomaly constraints~\eqref{eq:anomaly-linear}--\eqref{eq:anomaly-mixed}
      on the hypercharge weights $Y_R$.
      The explicit numerical coefficients of $a_4$ are collected in
      Appendix~\ref{app:seeley-dewitt}; only the structural
      form~\eqref{eq:topological-density} is required for the
      derivation of~\eqref{eq:anomaly-linear}--\eqref{eq:anomaly-mixed}.
    \end{proposition}

  \subsection{Hypercharge rigidity}
    \label{ssec:hypercharge-rigidity}

    \begin{proposition}[Hypercharge rigidity from spectral coherence]
      \label{prop:hypercharge-rigidity}
      Assume Hypotheses~\ref{hyp:spin-solder} and~\ref{hyp:weak-factor}, and that the projected spinorial
      sectors are precisely those of the admissible
      matter decomposition~\eqref{eq:matter-content}, with a supplied colour module
      $V_{\mathrm{color}}$ of dimension three (Section~\ref{ssec:colour-status}).
      That module is a hypothesis of this proposition, not a consequence of the preceding
      sections: the constraints below are weighted by $\dim(R)$, so the colour multiplicity
      enters the sums.
      Then $U(1)_Y$ gauge invariance of the projected spectral functional imposes
      the constraints~\eqref{eq:anomaly-linear}--\eqref{eq:anomaly-mixed} on the
      hypercharge weights $Y_R$.
      The allowed assignments are therefore not free parameters of the projected Dirac
      operator, but spectral coherence data constrained by $E_\Pi$.
    \end{proposition}

    \begin{proposition}[Minimal hypercharge assignment]
      \label{prop:minimal-hypercharge}
      If, in addition, the determinant line $L_Y = \wedge^2(E_{\mathrm{weak}})$ carries the minimal
      integral normalisation compatible with the $U(1)_Y$ action on the projected
      spinorial sectors, then the
      constraints~\eqref{eq:anomaly-linear}--\eqref{eq:anomaly-mixed} select the
      Standard Model hypercharge pattern up to an overall sign and rescaling.
    \end{proposition}

    The hypercharge weights are therefore not inserted as independent labels of matter
    fields.
    They are the weights of the determinant-line action that survive the spectral
    coherence condition imposed by the projected Dirac operator.
    The three-step derivation is:
    \begin{equation}
      \label{eq:hypercharge-chain}
      L_Y = \wedge^2(E_{\mathrm{weak}})
      \;\Rightarrow\;
      \text{support of } U(1)_Y
      \;\Rightarrow\;
      \delta_{\alpha Y} S_\Pi = 0
      \;\Rightarrow\;
      Y_R \text{ rigidity}.
    \end{equation}

  \section{Generation Multiplicity from the Rank-Three Admissible Invariant}
  \label{sec:generations}

  The preceding sections derive the spinorial electroweak bundle and the hypercharge
  consistency conditions, accounting for the leptonic spinorial sector.
  We now address the third structural output of the fermionic sector: generation
  multiplicity.
  The starting point is the admissible saturation invariant supplied as a selection rule
  by the spectral admissibility sub-programme~\cite{Beau2026a27} (O23 Theorem~3.1 proves
  the three-dimensionality of the neutral sector of a supplied spinor carrier; the
  identification of $\Sigma_c$ with that dimension is open),
  \begin{equation}
    \label{eq:sigma-n3}
    \sigma_c(n_3) = 3 .
  \end{equation}
  The purpose of the present section is to show that this invariant admits a spinorial
  multiplicity reading.
  No source identifies it with the spatial directions of the geometric branch: Q5b's effective
  co-metric has rank two on the spatial sector, with the central direction open
  (Q5b~\cite{Beau2026q5b}, Theorem~5.2 and open problem Q5b-O2), and the three-dimensional spatial
  sector of Q8 and Q11 carries a supplied spin-one action~\cite{BeauQ8,Beau2026q11}, not the O23
  invariant.

  \subsection{A geometric contrast and the spinorial reading}
    \label{ssec:two-functors}

    Let $\mathcal{A}_3$ denote the rank-three admissible sector selected at BI
    saturation:
    \begin{equation}
      \label{eq:admissible-sector}
      \dim \mathcal{A}_3 = 3 .
    \end{equation}
    Let $\{e_1, e_2, e_3\}$ be a basis of stable admissible directions in
    $\mathcal{A}_3$, orthogonal with respect to the admissibility form.
    The count three invites two readings of $\mathcal{A}_3$, of which only the second is
    constructed here.

    A geometric reading
    \begin{equation}
      \label{eq:geometric-functor}
      \mathcal{F}_{\mathrm{geom}} \colon \mathcal{A}_3 \longrightarrow
      \mathbb{R}^3_{\mathrm{space}}
    \end{equation}
    would send each stable admissible direction to a spatial tangent direction.
    It is recorded here only as a contrast: it is an interpretive parallel, not a construction
    supplied by the geometric branch, and nothing below uses it.

    The spinorial multiplicity reading is defined as follows.

    \begin{definition}[Spinorial multiplicity functor]
      \label{def:spinorial-functor}
      The spinorial multiplicity functor
      $\mathcal{F}_{\mathrm{spin}} \colon \mathcal{A}_3 \to \Mult(S_\Pi)$
      assigns to each stable admissible direction $e \in \mathcal{A}_3$ a copy
      $S_\Pi^{(e)}$ of the admissible spinor bundle, subject to the following conditions:
      \begin{enumerate}
        \item[\emph{(i)}]   $S_\Pi^{(e)} \cong S_\Pi$ as $\Spin(3,1)$ spinor bundles
        (Hypothesis~\ref{hyp:spin-solder});
        \item[\emph{(ii)}]  the gauge group $G_\Pi$ acts on $S_\Pi^{(e)}$ via the same
        representation as on $S_\Pi$, through the internal factors
        (Hypothesis~\ref{hyp:weak-factor});
        \item[\emph{(iii)}] the admissible projection $\Pi$ acts on $S_\Pi^{(e)}$ via
        the same kernel structure as on $S_\Pi$;
        \item[\emph{(iv)}]  distinct directions define orthogonal multiplicity sectors:
        the corresponding projectors $P_e$ satisfy
        \[
          P_e P_{e'} = 0 \quad (e \neq e'),
          \qquad
          [P_e, \rho(g)] = 0 \quad \text{for all } g \in G_\Pi .
        \]
      \end{enumerate}
    \end{definition}

    Condition~(iv) is the key constraint: it prevents the three copies from being
    interpreted as a gauge triplet.
    They form a multiplicity space on which $G_\Pi$ acts trivially, rather than a
    new non-trivial representation of $G_\Pi$.

  \subsection{Gauge-singlet character of the generation space}
    \label{ssec:gauge-singlet-generation}

    The generation space $\mathbb{C}^3_{\mathrm{gen}}$ is the indexing space of the
    functor $\mathcal{F}_{\mathrm{spin}}$, spanned by $\{e_1, e_2, e_3\}$.
    It is not the adjoint triplet $\operatorname{Sym}^2(S_\Pi)$, which carries the
    $\SU(2)$ adjoint action and belongs to the gauge bundle.
    By condition~(ii) of Definition~\ref{def:spinorial-functor}, the gauge group acts
    on each copy $S_\Pi^{(e_i)}$ in the same way as on $S_\Pi$; it acts trivially on
    the indexing labels $e_i$.
    Therefore
    \begin{equation}
      \label{eq:generation-singlet}
      \mathbb{C}^3_{\mathrm{gen}}
      \;\subset\;
      \ker\!\left(\operatorname{ad}_{\SU(2)} \oplus Y\right).
    \end{equation}

    \begin{theorem}[Generation multiplicity]
      \label{thm:generation-multiplicity}
      Assume the admissible saturation invariant $\sigma_c(n_3) = 3$~\cite{Beau2026a27}
      and the spinorial multiplicity functor $\mathcal{F}_{\mathrm{spin}}$ of
      Definition~\ref{def:spinorial-functor}.
      Then the fermionic matter bundle carries a three-dimensional gauge-singlet
      multiplicity factor:
      \begin{equation}
        \label{eq:fermionic-matter}
        S_\Pi^{\mathrm{fermion}}
        = S_\Pi \otimes \mathbb{C}^3_{\mathrm{gen}},
        \qquad
        \mathbb{C}^3_{\mathrm{gen}}
        \subset
        \ker\!\left(\operatorname{ad}_{\SU(2)} \oplus Y\right).
      \end{equation}
      The three-generation structure is the spinorial multiplicity image of the
      rank-three admissible invariant.
    \end{theorem}

    \begin{proof}
      By~\eqref{eq:sigma-n3}, $\mathcal{A}_3$ has dimension three and admits a basis of
      three mutually orthogonal stable admissible directions $\{e_1, e_2, e_3\}$.
      Applying $\mathcal{F}_{\mathrm{spin}}$ yields three copies
      $S_\Pi^{(e_1)}, S_\Pi^{(e_2)}, S_\Pi^{(e_3)}$.
      By condition~(iv) of Definition~\ref{def:spinorial-functor}, these copies define
      orthogonal projective sectors commuting with the $G_\Pi$-action, so their direct
      sum decomposes as
      \[
        \bigoplus_{i=1}^{3} S_\Pi^{(e_i)} \;\cong\; S_\Pi \otimes \mathbb{C}^3_{\mathrm{gen}}.
      \]
      The gauge group acts on each factor $S_\Pi^{(e_i)}$ via the spinorial
      representation and trivially on $\mathbb{C}^3_{\mathrm{gen}}$ by condition~(ii),
      giving~\eqref{eq:generation-singlet}.
    \end{proof}

    \begin{remark}
      \label{rem:space-generation-distinction}
      The equality $\dim \mathbb{R}^3_{\mathrm{space}} = \dim \mathbb{C}^3_{\mathrm{gen}} = 3$
      does not identify the two spaces.
      A spatial triplet would carry tangent structure, while the multiplicity image carries no
      electroweak charge; and no source supplies a geometric image of $\mathcal{A}_3$
      (Section~\ref{ssec:two-functors}).
      This distinction prevents the generation triplet from being confused with the
      adjoint $\SU(2)$ triplet $\operatorname{Sym}^2(S_\Pi)$.
    \end{remark}

  \section{Dynamic lifting of the static generation degeneracy}
  \label{sec:dyn-lift}

  Theorem~\ref{thm:generation-multiplicity} establishes the generation factor
  $\mathbb{C}^3_{\mathrm{gen}}$ as the gauge-singlet multiplicity space spanned by the three stable
  admissible directions $\{e_1, e_2, e_3\}$ of the rank-three sector $\mathcal{A}_3$
  (\eqref{eq:sigma-n3}, \eqref{eq:admissible-sector}), but leaves the quantitative
  inter-generation mass splitting open.
  This section isolates the structural mechanism by which such a splitting can be carried, prior to
  any amplitude analysis.
  The dimensionless generation structure is probed by the spectrum of $E_\Pi^2$
  (Definition~\ref{def:projective-endomorphism}) on $\mathbb{C}^3_{\mathrm{gen}}$.
  The physical mass scale and the amplitude of the splitting are not fixed at this stage; they are
  deferred to the normalisation of the admissibility cascade.
  The oriented diagnostic of this cascade splitting is developed in~\cite{Beau2026q11of}; its chiral weight is
  identified there and in~\cite{Beau2026cho} as a Weyl representation type rather than a function of the cascade
  generator, with the frontier normalisation fixed maximally in the inherited central-phase lift.

  We use a supplied model of the internal structure of the rank-three sector.
  As supplied model data in O30~\cite{Beau2026a34}, the admissible sector is the symmetric square
  \begin{equation}
    \label{eq:a3-sym2}
    \mathcal{A}_3 \;\cong\; \Sym^2(V)
  \end{equation}
  of the two-dimensional admissible carrier $V$, the supplied spinor carrier of
  O23~\cite{Beau2026a27} (Section~4.1), carrying the weight grading of the Born--Infeld datum.
  The orthogonal admissible basis $\{e_1, e_2, e_3\}$ is thereby reorganised into the weight basis
  $\{e_0, e_+, e_-\}$ of $J_3$-weights $\{0, +1, -1\}$, with $e_0$ the central weight-zero
  direction.
  This identification is supplied, not derived: the finite-precision diagnostics of
  O29~\cite{Beau2026a33} are compatible with it and do not identify it.
  It probes the internal structure of $\mathbb{C}^3_{\mathrm{gen}}$ and does not alter its gauge-singlet
  character (Theorem~\ref{thm:generation-multiplicity}).
  The reality structure is the internal antilinear parity $J_\Pi = HK$ of
  Subsection~\ref{ssec:proof-bi-left}, realising the parity $c \leftrightarrow q-c$, which fixes $e_0$
  and exchanges $e_+ \leftrightarrow e_-$ antilinearly.
  It uses the supplied Born--Infeld datum and Lemma~\ref{lem:bi-grading-anticommute}, not
  Hypothesis~\ref{hyp:spin-solder}.

  \subsection{Static obstruction}
    \label{ssec:dg-static}

    \begin{proposition}[Static obstruction]
      \label{prop:dg-static}
      Let $C \in \End(\mathbb{C}^3_{\mathrm{gen}})$ be a static restriction of $E_\Pi^2$ to the
      generation factor satisfying
      \begin{enumerate}
        \item[\emph{(i)}]  $J_\Pi$-reality, $J_\Pi\, C\, J_\Pi^{-1} = C^{\sharp}$, and
        \item[\emph{(ii)}] weight-sector preservation, $[C, J_3] = 0$.
      \end{enumerate}
      Then $C_{++} = C_{--}$: the outer pair $\{e_+, e_-\}$ is degenerate, and $C$ admits no real
      splitting of the two light generations.
    \end{proposition}

    \begin{proof}
      Weight preservation~(ii) makes $C$ diagonal in the weight basis, so its doublet block is
      $\operatorname{diag}(C_{++}, C_{--})$.
      For the real representative, $J_\Pi$-reality~(i) reads $P\, C\, P = C$, where
      $P = \left(\begin{smallmatrix} 1 & 0 & 0 \\ 0 & 0 & 1 \\ 0 & 1 & 0 \end{smallmatrix}\right)$
      exchanges $e_+ \leftrightarrow e_-$ and fixes $e_0$.
      Conjugating the diagonal doublet block by
      $P|_{\{e_+,e_-\}} = \left(\begin{smallmatrix} 0 & 1 \\ 1 & 0 \end{smallmatrix}\right)$ swaps
      its two entries, so $P\,C\,P = C$ forces $C_{++} = C_{--}$.
    \end{proof}

    \begin{remark}[Sharpness of the hypotheses]
      \label{rem:dg-sharpness}
      $J_\Pi$-reality alone is not sufficient.
      The condition $P M P = M$ on the doublet admits
      $M = \left(\begin{smallmatrix} a & b \\ b & a \end{smallmatrix}\right)$, whose eigenvalues
      $a \pm b$ are split whenever $b \neq 0$: the parity fixes the diagonal entries equal but does
      not diagonalise the block.
      The degeneracy requires the static weight-sector preservation~(ii) in addition to~(i).
      The matrix $\operatorname{diag}(1, \tfrac12, \tfrac12)$ satisfies both hypotheses; up to
      normalisation and ordering of the weight basis it is the form of the conditional covariance of
      O30~\cite{Beau2026a34}, whose identification with a restriction of $E_\Pi^2$ is not
      available~\cite{Beau2026prs}.
    \end{remark}

  \subsection{The oriented odd sector}
    \label{ssec:dg-odd}

    \begin{proposition}[Oriented odd sector]
      \label{prop:dg-odd}
      Let $\mathcal{T}_\tau \in \End(\mathbb{C}^3_{\mathrm{gen}})$ be a lift required to
      \begin{enumerate}
        \item[\emph{(i)}]   annihilate the central level, $\mathcal{T}_\tau e_0 = 0$,
        \item[\emph{(ii)}]  preserve the outer doublet,
        $\mathcal{T}_\tau(\{e_+, e_-\}) \subseteq \{e_+, e_-\}$, and
        \item[\emph{(iii)}] be $J_\Pi$-odd, $J_\Pi\, \mathcal{T}_\tau\, J_\Pi^{-1} = -\mathcal{T}_\tau$.
      \end{enumerate}
      Then the admissible lift space is two-dimensional, spanned by the weight generator
      $J_3 = \operatorname{diag}(0, 1, -1)$ and the doublet rotation
      $R_{\mathrm{mix}} = \left(\begin{smallmatrix} 0 & 0 & 0 \\ 0 & 0 & 1 \\ 0 & -1 & 0
      \end{smallmatrix}\right)$.
      The direction $J_3$ is the unique non-mixing direction producing a real mass split.
      The Carnot dilation $D = \operatorname{diag}(2, 1, 1)$ is $J_\Pi$-even and therefore excluded
      from this sector.
    \end{proposition}

    \begin{proof}
      Conditions~(i)--(ii) confine $\mathcal{T}_\tau$ to a real $2 \times 2$ doublet block $M$.
      Condition~(iii) for the real representative is $P_2\, M\, P_2 = -M$ with
      $P_2 = \left(\begin{smallmatrix} 0 & 1 \\ 1 & 0 \end{smallmatrix}\right)$, giving
      $M = \left(\begin{smallmatrix} a & b \\ -b & -a \end{smallmatrix}\right)$, a two-parameter
      traceless family with $\lambda^2 = a^2 - b^2$.
      The $a$-direction is $J_3|_{\mathrm{doublet}}$, with real eigenvalues $\pm a$ and no mixing;
      the $b$-direction is $R_{\mathrm{mix}}$, with eigenvalues $\pm i b$ at $a = 0$ and pure
      mixing.
      A real split exists for any $|a| > |b|$, but $J_3$ ($b = 0$) is the only non-mixing splitting
      direction.
      The Carnot dilation satisfies $P D P = D$, hence even, and is excluded.
    \end{proof}

    \begin{remark}[Mixing direction]
      \label{rem:dg-mixing}
      The odd sector contains both a real-split direction and a mixing direction.
      The real cascade model of Subsection~\ref{ssec:dg-bridge} activates the former but not the
      latter; the latter is expected, if present, to require the complex metaplectic phase.
    \end{remark}

  \subsection{The cascade bridge}
    \label{ssec:dg-bridge}

    Let $g = S_X(t)\, S_Y(s)$ be the oriented metaplectic step, the ordered product of the two
    horizontal Heisenberg generators in $\SL(2)$ form, and let
    \[
      \mathcal{G}_g := \log g ,
    \]
    lifted to $\Sym^2(V)$ and transported to $\mathbb{C}^3_{\mathrm{gen}}$ by~\eqref{eq:a3-sym2}, be
    its per-step generator.
    We call $\mathcal{G}_g$ the oriented cascade generator; the name designates this model object
    and nothing else.
    Its projection onto the $J_\Pi$-odd sector of Proposition~\ref{prop:dg-odd} is computed below.

    Identifying $\mathcal{G}_g$ with an emergent ordering derivative $\partial_\tau$ would be a
    separate hypothesis, and no source supplies it.
    The ordering derivative of~\cite{Beau2026q11} is defined only on functions of the ordering
    parameter $\tau = n\,\Delta\tau$, with a free step $\Delta\tau$, and defines no operator on a
    representation space of $\SU(2)$ (Q11 Modelling identification~3.1 and Remark~3.2).
    Nothing in this section uses that identification: the orientation of $g$ is the order of the
    product, and the results below are statements about $\mathcal{G}_g$.

    \begin{proposition}[Cascade bridge]
      \label{prop:dg-bridge}
      In this model the $J_\Pi$-odd projection of the oriented cascade generator
      $\mathcal{G}_g$ on $\mathbb{C}^3_{\mathrm{gen}}$ is
      \begin{equation}
        \label{eq:dg-projection}
        \Pi_{J_\Pi\text{-odd}}\!\left( \mathcal{G}_g\big|_{\mathbb{C}^3_{\mathrm{gen}}} \right)
        = \alpha\, J_3 + \mu\, R_{\mathrm{mix}}, \qquad \alpha \neq 0 .
      \end{equation}
      The real-split coefficient $\alpha$ is non-vanishing already for symmetric shears $t = s$,
      where it is produced by the order of the product, that is by the non-commutativity
      $[X, Y] = Z$.
      Reversal of the cascade order, $g = S_Y(t)\, S_X(s)$, reverses the sign,
      $\alpha_{\mathrm{rev}} = -\alpha$.
      The mixing coefficient $\mu$ vanishes throughout the real symplectic model.
    \end{proposition}

    A supporting diagnostic computation builds the rank-three factor from $V = \mathbb{C}^2$ by $\Sym^2$,
    takes $\log(g)$ on the lifted representation, and projects its $J_\Pi$-odd part onto $J_3$ and
    $R_{\mathrm{mix}}$ in the Hilbert--Schmidt inner product.
    It returns $\alpha = +0.156$ at $t = s = 0.40$, with $\alpha_{\mathrm{rev}} = -0.156$ and
    $\mu = 0$.

  \subsection{The exit-deficit dictionary}
    \label{ssec:dg-dictionary}

    The amplitude target of the lift (Subsection~\ref{ssec:dg-bridge}) is defined not by the raw
    eigenvalues of $E_\Pi^2$ but by their \emph{deficit} relative to the central level.
    This is the observable that matches the spectral admissibility hierarchy, where the relevant
    quantity is the support-exit distance from the non-exiting central ADE level
    $\lambda_2 = 1$ of O1~\cite{Beau2026a6}, whose levels $\{\tfrac56, 1, \tfrac54\}$ are the
    non-neutral ord-5 spectral set of~\cite{Beau2026a4} normalised by its central value; the choice of
    that generator set over the neutral ord-4 set is open there.

    \begin{definition}[Exit deficit]
      \label{def:dg-exit-deficit}
      For $s_i \in \mathrm{Spec}(E_\Pi^2|_{\mathbb{C}^3_{\mathrm{gen}}})$ with $s_0$ the central
      ($J_\Pi$-fixed) eigenvalue, the exit deficit of $e_i$ is
      \[
        d_{\mathrm{exit}}(e_i) := s_0 - s_i .
      \]
    \end{definition}

    \begin{proposition}[Spectrum--ADE dictionary]
      \label{prop:dg-dictionary}
      For the dynamical spectrum
      $\mathrm{Spec}(E_\Pi^2|_{\mathbb{C}^3_{\mathrm{gen}}}) = \{s_0 = 1,\ s_\pm = \tfrac12 \pm
      \varepsilon\}$ the exit deficits are
      \[
        d_{\mathrm{exit}}(e_0) = 0, \qquad d_{\mathrm{exit}}(e_\pm) = \tfrac12 \mp \varepsilon .
      \]
      The central direction $e_0$, fixed by $J_\Pi$, maps to the non-exiting level $\lambda_2 = 1$.
      The two outer deficits map to the exiting ADE distances
      $\{|\lambda_1 - 1|, |\lambda_3 - 1|\} = \{\tfrac16, \tfrac14\}$ up to a single common scale
      $\kappa$, and the external ratio condition
      \[
        \frac{d_{\mathrm{exit}}(e_-)}{d_{\mathrm{exit}}(e_+)}
        = \frac{\tfrac12 + \varepsilon}{\tfrac12 - \varepsilon}
        = \frac{|\lambda_3 - 1|}{|\lambda_1 - 1|}
        = \frac32
        \qquad\Longrightarrow\qquad
        \varepsilon = \tfrac{1}{10}, \quad \kappa = \tfrac{5}{12} .
      \]
      A larger deficit corresponds to an earlier support exit, so
      $d_{\mathrm{exit}}(e_0) = 0 < d_{\mathrm{exit}}(e_+) < d_{\mathrm{exit}}(e_-)$ reproduces the
      exit-order mass ordering $M_3 > M_1 > M_2$ of O1~\cite{Beau2026a6} on the levels.
      Reading it as an ordering of the three generations requires the level-to-generation map, which
      is not established~\cite{Beau2026a6}.
    \end{proposition}

    \begin{proof}
      The deficits follow from Definition~\ref{def:dg-exit-deficit} applied to the spectrum.
      With $\varepsilon = \tfrac1{10}$ the outer deficits are $\{\tfrac25, \tfrac35\}$, in ratio
      $\tfrac32$, matching $\{|\lambda_1 - 1|, |\lambda_3 - 1|\} = \{\tfrac16, \tfrac14\}$ also in
      ratio $\tfrac32$; the single scale $\kappa$ solving $\kappa\, d_{\mathrm{exit}} = |\lambda - 1|$
      is $\kappa = (\tfrac14)/(\tfrac35) = (\tfrac16)/(\tfrac25) = \tfrac{5}{12}$, consistent on both
      levels.
    \end{proof}

    \begin{remark}[Resolution of the apparent inversion]
      \label{rem:dg-inversion}
      The raw eigenvalues $s_i$ measure the static coherence carried by the generation factor, for
      which the central direction $e_0$ is maximal; the exit deficits $d_{\mathrm{exit}}$ measure the
      deficit of exit relative to the central level, for which $e_0$ is zero (non-exiting).
      The two observables are linked by $d_{\mathrm{exit}} = s_0 - s$, so the apparent conflict
      between the central-maximal spectrum here and the central-midpoint level $\lambda_2 = 1$ of the
      ADE picture dissolves: they live in different observables related by a deficit.
      The value $\varepsilon = \tfrac1{10}$ is fixed by imposing the external ADE ratio on the exit
      deficits: it is a matching value, not an independent prediction, and it acquires predictive
      content only if $\varepsilon$ is computed from the cascade normalisation without that ratio.
    \end{remark}

    The common scale $\kappa = \tfrac{5}{12}$ is the passage between the normalised deficit here and
    the ADE distance; it is not yet a mass.
    A mass hierarchy would require an amplification law; the prescription
    $\beta^* = 1/(\delta_{\mathrm{pair}} + \tfrac12) \approx 0.126$ is phenomenological and not
    derived~\cite{Beau2026a26}, so the hierarchy is deferred together with the cascade normalisation of
    Subsection~\ref{ssec:dg-status}.

  \subsection{Status and the amplitude problem}
    \label{ssec:dg-status}

    The qualitative segment is closed:
    \begin{equation}
      \label{eq:dg-chain}
      \text{static obstruction} \;\longrightarrow\; \text{oriented lift} \;\longrightarrow\;
      \mathcal{G}_g \text{ carries } \alpha J_3, \quad \alpha \neq 0 .
    \end{equation}
    The splitting of the two light generations is forbidden statically by $J_\Pi$-reality together
    with weight preservation (Proposition~\ref{prop:dg-static}), lives in the two-dimensional
    $J_\Pi$-odd sector (Proposition~\ref{prop:dg-odd}), and is carried, in the metaplectic step
    model of Subsection~\ref{ssec:dg-bridge}, by the oriented cascade generator $\mathcal{G}_g$
    through its non-zero $J_3$ projection (Proposition~\ref{prop:dg-bridge}).

    The value of $\alpha$ obtained in the real shear model is not a mass prediction.
    Only its non-vanishing and its sign reversal under cascade reversal are structural.
    The amplitude problem is deferred to the normalisation of the admissibility cascade.
    Concretely, the exit-deficit ratio of Proposition~\ref{prop:dg-dictionary} enters through a
    normalised coefficient $\varepsilon$, not through $\alpha$ itself:
    \begin{equation}
      \label{eq:dg-amplitude}
      \frac{\tfrac12 + \varepsilon}{\tfrac12 - \varepsilon} = \frac32, \qquad
      \varepsilon = \mathcal{N}_{\mathrm{casc}}\, |\alpha| .
    \end{equation}
    The diagnostic architecture test further separates the two roles entering
    \eqref{eq:dg-amplitude}.
    The $J_3$ coefficient is an angular, metaplectic contribution: it belongs to the traceless
    adjoint sector of $\End(\Sym^2 V)$ and preserves the Hilbert--Schmidt norm.
    By contrast, the decay of $\sigma_{\mathrm{pair}}(n)$ is a radial capacity effect: it changes
    the projected norm and therefore belongs to the scalar sector supplied by the non-injective
    projection.
    Equivalently,
    \[
      \langle \mathbf{1}, J_3 \rangle_{\mathrm{HS}} = 0,
    \]
    so the generation split and the capacity decay are distinct, Hilbert--Schmidt orthogonal
    generators.

    Thus the amplitude architecture is factorised rather than unified.
    The same cascade arrow orients both observables: reversing the metaplectic order reverses the
    sign of the $J_3$ split, while the cumulative projected capacity
    $I(n) = \sigma_{\mathrm{pair}}(0) - \sigma_{\mathrm{pair}}(n)$ fixes the monotone orientation
    of the radial projection sector.
    However, the two generators are not identical.
    Consequently $\mathcal{N}_{\mathrm{casc}}$ in \eqref{eq:dg-amplitude} should be read as a
    transfer constant between the angular metaplectic split and the radial projective decay, not as
    a parameter-free function of $\delta_{\mathrm{pair}}$ alone.
    The remaining amplitude problem is therefore to compute this transfer constant from the coupling
    of the two cascade observables, rather than to identify $\alpha$ with the capacity exponent.

    The transfer observable must therefore be defined with respect to the normalised cumulative
    projected capacity
    \[
      \widehat{I}(n)
      =
      \frac{\sigma_{\mathrm{pair}}(0) - \sigma_{\mathrm{pair}}(n)}
      {\sigma_{\mathrm{pair}}(0) - \sigma_{\mathrm{pair}}(n_{\mathrm{sat}})} ,
      \qquad
      \widehat{I}(n_{\mathrm{sat}}) = 1 .
    \]
    The corresponding angular--radial transfer profile is
    \[
      \varepsilon_n
      =
      \sum_{m \le n} \mathcal{K}_{\mathrm{ar}}(m)\, \Delta\widehat{I}(m),
      \qquad
      \Delta\widehat{I}(m) = \widehat{I}(m) - \widehat{I}(m-1) .
    \]
    The structurally admissible locking law is the one in which the angular split is accumulated per
    unit of projected capacity,
    \[
      \Delta\varepsilon(m) = \mathcal{K}_{\mathrm{ar}}\, \Delta\widehat{I}(m) .
    \]
    It is selected by the factorised architecture above: the metaplectic $J_3$ component does not
    itself carry radial decay, while the radial projection sector fixes the common orientation
    through $\widehat{I}(n)$.
    Coupling the split instead to the bare cascade step would make the accumulated amplitude depend
    on the parametrisation depth, and coupling it to the residual capacity would tie the split to the
    unprojected rather than to the realised part of the cascade.
    Both alternatives are therefore diagnostics, not the transfer law.

    Under this locking law the transfer profile is constant:
    \[
      \mathcal{K}_{\mathrm{ar}}(m) = \mathcal{K}_{\mathrm{ar}},
      \qquad
      \varepsilon_{\mathrm{sat}}
      =
      \sum_{m \le n_{\mathrm{sat}}} \mathcal{K}_{\mathrm{ar}}\, \Delta\widehat{I}(m)
      =
      \mathcal{K}_{\mathrm{ar}} .
    \]
    Thus the cascade normalisation in \eqref{eq:dg-amplitude} is a pure angular--radial transfer
    constant.
    This fixes the profile and the saturation status of the amplitude, but not its absolute value.
    The latter still requires the physical normalisation of the metaplectic $J_3$ generator; the
    shear value $\alpha \simeq 0.156$ remains a diagnostic coefficient, not an input to the mass
    prediction.


  \section{The Colour Sector}
  \label{sec:colour}

  The preceding sections derive the spinorial electroweak structure and the
  generation multiplicity factor, accounting for the leptonic spinorial sector.
  The quark sector requires, in addition, a colour representation.
  In the present framework, the colour factor is not introduced by enlarging the
  spinorial construction itself, but by tensoring the admissible spinor bundle with a colour
  representation carried as a supplied input.
  The gauge-fibre programme does not derive one: see Section~\ref{ssec:colour-status}.

  Let $V_{\mathrm{color}}$ denote a supplied three-dimensional colour representation; no
  admissible colour sector is derived (Section~\ref{ssec:colour-status}).
  The quark spinor bundle is defined as
  \begin{equation}
    \label{eq:quark-bundle}
    S_\Pi^{\mathrm{quark}} = S_\Pi \otimes V_{\mathrm{color}} .
  \end{equation}
  The full fermionic matter bundle targeted by the construction is therefore
  \begin{equation}
    \label{eq:matter-bundle-full}
    S_\Pi^{\mathrm{matter}}
    =
    \left(
      S_\Pi \oplus (S_\Pi \otimes V_{\mathrm{color}})
    \right)
    \otimes \mathbb{C}^3_{\mathrm{gen}} .
  \end{equation}

  \subsection{Status of the colour factor}
    \label{ssec:colour-status}

    The construction of $S_\Pi$, $L_Y$, and $\mathbb{C}^3_{\mathrm{gen}}$ does not
    depend on the colour hypothesis.
    The factor $V_{\mathrm{color}}$ does, and it is a supplied input.

    No identification of $\SU(3)$ as a colour gauge group is available.
    O31~\cite{Beau2026a35} is a withdrawal notice.
    Among the results it withdraws are the co-admissibility of the individual capacity profiles
    $\sigma_{c_1}$, $\sigma_{c_2}$, $\sigma_{c_3}$ --- which is the content of
    $[H\text{-color}]_{\mathrm{pointwise}}$ and of its Proposition~4.23 --- the uniqueness of
    $\SU(3)$ among connected compact subgroups acting irreducibly on $\mathbb{C}^3$; the
    irreducibility of the group of fibre-preserving transformations of the colour eigenspace; the
    arithmetic criterion $q \equiv 1 \pmod 3$ for the existence of colour triplets; and every
    derivation of a gauge factor.
    Three counterexamples are given there: $\operatorname{SO}(3) \subset \SU(3)$, a proper
    connected compact subgroup acting irreducibly on $\mathbb{C}^3$ with determinant one; the diagonal torus $T^2$,
    which satisfies every constraint that was imposed while acting reducibly; and
    $\{1, 2, q-3\}$, an admissible triplet for every prime $q \geq 7$ irrespective of
    $q \bmod 3$.
    Nothing withdrawn there is replaced by a weaker positive statement.

    O32~\cite{Beau2026a36} measures the finite-$q$ residual variance of the capacity profiles.
    Those are measurements; this paper takes no position on their interpretation and does not read
    them as identifying a group.

    $V_{\mathrm{color}} \in \operatorname{Rep}(\SU(3))$ of dimension three is therefore carried
    as a hypothesis, and the quark sector~\eqref{eq:quark-bundle} is a conditional output of the
    present paper.
    Two results of this paper rest on it: the quark carrier itself, and the hypercharge selection.
    Proposition~\ref{prop:hypercharge-rigidity} carries the supplied module in its hypothesis, and
    Proposition~\ref{prop:minimal-hypercharge}, which is where the selection of the Standard Model
    pattern is stated, inherits it through its ``If, in addition'' clause: the cubic
    constraint~\eqref{eq:anomaly-cubic} vanishes on that pattern precisely at colour multiplicity
    three.
    The $V{-}A$ structure, the spinorial lift of BI parity, the leptonic carrier and the generation
    factor $\mathbb{C}^3_{\mathrm{gen}}$ do not.

    \begin{proposition}[Colour sector extension]
      \label{prop:colour-extension}
      Let $V_{\mathrm{color}}$ be a supplied three-dimensional representation of a colour group,
      carried as a hypothesis rather than derived (Section~\ref{ssec:colour-status}).
      Then the spinorial matter bundle admits the colour-coupled extension
      \begin{equation}
        \label{eq:colour-extension}
        S_\Pi^{\mathrm{matter}}
        =
        \left(
          S_\Pi \oplus (S_\Pi \otimes V_{\mathrm{color}})
        \right)
        \otimes \mathbb{C}^3_{\mathrm{gen}} .
      \end{equation}
      The first summand is the leptonic sector.
      The second summand is the quark sector.
      The generation factor $\mathbb{C}^3_{\mathrm{gen}}$ remains a gauge-singlet
      multiplicity space under both summands.
    \end{proposition}

    \begin{remark}
      \label{rem:colour-does-not-affect-spinorial-construction}
      The colour extension does not modify the construction of the admissible spinor
      bundle $S_\Pi$.
      It adds an independent gauge representation factor, which is supplied rather than derived.
      The logical status of the fermionic construction is therefore stratified:
      \begin{equation}
        \label{eq:status-table}
        \begin{cases}
          S_\Pi \otimes \mathbb{C}^3_{\mathrm{gen}}
          & \text{electroweak/leptonic carrier, no colour condition,} \\[4pt]
          S_\Pi \otimes V_{\mathrm{color}} \otimes \mathbb{C}^3_{\mathrm{gen}}
          & \text{quark carrier, conditional on the supplied } V_{\mathrm{color}}.
        \end{cases}
      \end{equation}
      Both rows are stated relative to the supplied Heisenberg carrier of
      Subsection~\ref{ssec:weil-metaplectic} and to Hypotheses~\ref{hyp:spin-solder}
      and~\ref{hyp:weak-factor}, which the whole physical reading inherits; the
      distinction drawn here is the additional colour condition carried by the second row alone.
    \end{remark}

    \newpage
  \section{Conclusion and Open Directions}
  \label{sec:conclusion}

  This paper states what the fermionic sector of the spectral approach requires beyond the admissible
  Weil fibre.
  Fermions are not introduced as external matter fields added to the admissible bundle; they are read
  from the Weil module carried by the admissible fibre on a supplied Heisenberg carrier, through the
  complexification of its metaplectic Lie algebra.

  The structural chain of the paper is
  \[
    \Pi
    \Rightarrow
    F_n \simeq V_\rho
    \Rightarrow
    \Mp(2,\mathbb{R})
    \Rightarrow
    \mathfrak{mp}(2,\mathbb{R})_{\mathbb{C}} \simeq \mathfrak{sl}_2(\mathbb{C})
    \overset{\text{[H-Spin]}}{\leadsto}
    \Spin(3,1) \simeq \SL(2,\mathbb{C})
    \Rightarrow
    S_\Pi .
  \]
  Every arrow before the one marked [H-Spin] is proved relative to the supplied carrier.
  On the abstract doublet, the symmetric square is the adjoint of $\mathfrak{sl}_2(\mathbb{C})$ and the
  exterior square is the trivial line.
  The arrow marked [H-Spin] is Hypothesis~\ref{hyp:spin-solder}: a four-dimensional Lorentzian co-metric
  on the effective base with a spin structure whose structure group is the $\SL(2,\mathbb{C})$ acting on
  the doublet.
  Under it, $\operatorname{Sym}^2(S_L)$ is the Lorentz sector $\operatorname{ad}(P_\Pi)$ and
  $\wedge^2(S_L)$ is trivial.
  The electroweak data come from a second factor: by Schur's lemma a weak action commuting with Lorentz
  transformations needs a distinct rank-two bundle $E_{\mathrm{weak}}$, Hypothesis~\ref{hyp:weak-factor},
  on which
  \[
    \operatorname{Sym}^2(E_{\mathrm{weak}})
    \Rightarrow
    \operatorname{ad}_{\mathbb{C}}(\mathfrak{su}(2)_L),
    \qquad
    \wedge^2(E_{\mathrm{weak}})
    \Rightarrow
    L_Y ,
  \]
  the determinant line carrying $U(1)_Y$ because the structure group of $E_{\mathrm{weak}}$ is $U(2)$.
  The admissible doublet can supply at most one of the two factors, and this paper does not construct
  the other.

  Under Hypothesis~\ref{hyp:spin-solder}, the projected Dirac operator
  \[
    \mathcal{D}_{\Pi,g,A}
    =
    \Pi_S \circ \mathcal{D}_{g,A} \circ \Pi_S^*
  \]
  has, for a symbol-compatible lift, a Laplace-type square with standard Clifford principal symbol and a
  projective zero-order endomorphism $E_\Pi$:
  \[
    \mathcal{D}_{\Pi,g,A}^2
    =
    -(\nabla^{S,A})^2
    +
    \frac{R}{4}
    +
    \frac{1}{2} \gamma^\mu\gamma^\nu F^a_{\mu\nu}T^a_R
    +
    E_\Pi .
  \]
  The endomorphism $E_\Pi$ is the internal spectral residue of non-injective projection in the spinorial
  sector.
  In this sense, the analogue of the finite internal Dirac operator is not postulated as additional data:
  it is induced by the projective construction.

  The chiral statements separate into two kinds.
  The first concerns the Lorentz chirality: with the Born--Infeld datum supplied by (H1)--(H2) of O30,
  the internal parity $J_\Pi = HK$ satisfies $J_\Pi^2 = -1$, and under Hypothesis~\ref{hyp:spin-solder}
  it is the unique covariant antilinear lift and reverses chirality
  (Theorem~\ref{thm:spinorial-bi-lift}); on the orientation-compatible branch the projective residue is
  left-admissible,
  \[
    P_R E_\Pi P_R = 0 .
  \]
  The second concerns the weak interaction: its chiral selection, the $V{-}A$ structure, is the
  left-handed doublet assignment of Hypothesis~\ref{hyp:weak-factor}, which left-admissibility is
  compatible with but does not select (Remark~\ref{rem:lorentz-vs-weak-chirality}).
  Under both hypotheses, $U(1)_Y$ invariance of the projected spectral functional imposes the local
  anomaly-cancellation condition
  \[
    \operatorname{Tr}_{S_\Pi}
    \!\left(
        \gamma_5\,Y\,\mathcal{A}_\Pi(x)
    \right)
    =
    0
    \qquad
    \text{for all } x\in M_{\mathrm{eff}} ,
  \]
  where $\mathcal{A}_\Pi$ is the $\gamma_5$-weighted $a_4$ Seeley--DeWitt density
  of $\mathcal{D}_{\Pi,g,A}^2$.
  Anomaly cancellation then sits at the same $a_4$ level that the bosonic synthesis reads as the
  Yang--Mills response.

  The third output is generation multiplicity, conditional on the supplied rank-three
  selection rule of O23.
  The spinorial multiplicity functor sends the rank-three admissible invariant $\sigma_c(n_3)=3$ to a
  gauge-singlet generation factor:
  \[
    \mathbb{C}^3_{\mathrm{gen}}
    \subset
    \ker\!\left(\operatorname{ad}_{\SU(2)}\oplus Y\right).
  \]
  This generation space is neither the adjoint triplet nor a spatial triplet; no source identifies the
  O23 invariant with spatial directions.
  It is a multiplicity space commuting with the gauge action.
  The resulting fermionic carrier is
  \[
    S_\Pi^{\mathrm{fermion}}
    =
    S_\Pi \otimes \mathbb{C}^3_{\mathrm{gen}}
  \]
  in the electroweak/leptonic sector, with the weak factor of Hypothesis~\ref{hyp:weak-factor}.

  The colour-coupled quark sector is obtained by tensoring with the colour module:
  \[
    S_\Pi^{\mathrm{quark}}
    =
    S_\Pi \otimes V_{\mathrm{color}} .
  \]
  This step is conditional: $V_{\mathrm{color}}$ is a supplied three-dimensional representation,
  and O31~\cite{Beau2026a35} withdraws every derivation of one
  (Section~\ref{ssec:colour-status}).
  Accordingly, the matter carrier addressed by the construction is
  \[
    S_\Pi^{\mathrm{matter}}
    =
    \left(
      S_\Pi \oplus (S_\Pi\otimes V_{\mathrm{color}})
    \right)
    \otimes
    \mathbb{C}^3_{\mathrm{gen}},
  \]
  whose first summand carries the supplied Heisenberg carrier and the two hypotheses, and whose second
  carries the supplied colour module as well.

  Several open directions are sharply delimited.
  First, the two hypotheses themselves: constructing the solder of Hypothesis~\ref{hyp:spin-solder}
  from the metaplectic action, and a second rank-two factor for Hypothesis~\ref{hyp:weak-factor}, or
  showing that the admissible fibre supplies neither.
  Second, the minimal integral normalisation of the determinant line $L_Y$ must be
  connected explicitly to the Standard Model hypercharge pattern.
  The anomaly constraints restrict the allowed weights, while the determinant-line
  normalisation fixes their scale.
  Third, the amplitude of the generation splitting remains open: Section~\ref{sec:dyn-lift}
  fixes the qualitative mechanism in the metaplectic step model, but the coefficient $\varepsilon$ must
  still be derived from the cascade normalisation rather than matched.
  Using the O12 fingerprint~\cite{Beau2026a16}, whose pure Weil frequency is $B_c=\sum_i c_i b_i$, the
  residual reflection $R_b$ left by the oriented-frontier diagnostic is forced to be the central Weil
  lift $W(-I)=\mathrm{FT}^2$: preserving $A_c$, including its cocycle term, sends
  $(a_i,b_i)\mapsto(-a_i,-b_i)$.
  As the centre of $\SL(2,\mathbb{C})$, $R_b$ is scalar on $S_L\oplus S_R$ under
  Hypothesis~\ref{hyp:spin-solder}, so $[\gamma_5,R_b]=0$, and the chirality grading removes the residual
  obstruction without a spacetime parity.
  Fourth, possible generation mixing requires the complex metaplectic phase, since the real
  symplectic cascade model activates the $J_3$ split direction but not the $R_{\mathrm{mix}}$
  direction.

  A further direction concerns the physical content of $E_\Pi$.
  Since $E_\Pi$ is the internal spectral residue of non-injective projection, it is
  the natural place where fermion masses, effective Yukawa couplings, and
  inter-generation mixing should enter.
  Section~\ref{sec:dyn-lift} takes the first step on this direction: the
  inter-generation splitting is statically obstructed on $\mathbb{C}^3_{\mathrm{gen}}$ by
  $J_\Pi$-reality together with weight preservation (Proposition~\ref{prop:dg-static}), lives in a
  two-dimensional $J_\Pi$-odd lift sector (Proposition~\ref{prop:dg-odd}), and is carried in the
  metaplectic step model by the oriented cascade generator $\mathcal{G}_g$ through its non-zero $J_3$
  projection (Proposition~\ref{prop:dg-bridge}).
  Identifying $\mathcal{G}_g$ with an emergent ordering derivative is not supplied.
  The amplitude and the comparison with the charged-lepton mass hierarchy~\cite{Beau2026a4,Beau2026a29}
  are among the open directions identified above.

  \subsection*{Interpretive outlook}

  The following is a reading of the results, not a result.
  The Standard Model does not treat spin and weak isospin as one structure: a left-handed electron is a
  Lorentz spinor and, separately, a member of a weak doublet.
  The present analysis reaches the same separation from the other side.
  The admissible Weil fibre supplies one doublet and its algebra, and that algebra is exactly the
  algebra of both spin and weak isospin; but one copy cannot carry both roles.
  Where the fermionic sector meets spacetime and where it meets the weak interaction are therefore two
  distinct questions, and the question of why weak interactions are left-handed becomes the question of
  how two such factors are joined.
  The remaining problems are specific: two hypotheses on how the fibre is soldered, and rigidity problems
  for $E_\Pi$ and $L_Y$.

  \newpage
  \begin{appendix}

    \section{Seeley--DeWitt Coefficients for the Projected Dirac Square}
      \label{app:seeley-dewitt}

      Under symbol-compatibility (Section~\ref{ssec:projective-correction}),
      $\mathcal{D}_{\Pi,g,A}^2$ is a Laplace-type operator of the form
      $D = -(\nabla^{S,A})^2 + E$ with $E$ as in~\eqref{eq:endomorphism-E}.
      The standard Seeley--DeWitt expansion gives~\cite{Vassilevich2003}
      \begin{align}
        \label{eq:sdw-a0}
        a_0(x, D) &= (4\pi)^{-2} \operatorname{tr}(\mathrm{id}), \\
        \label{eq:sdw-a2}
        a_2(x, D) &= (4\pi)^{-2} \operatorname{tr}\!\left(\tfrac{R}{6}\,\mathrm{id} - E\right), \\
        \label{eq:sdw-a4}
        a_4(x, D) &= (4\pi)^{-2} \operatorname{tr}\!\left(
                                                      \tfrac{1}{2}E^2 - \tfrac{1}{6}RE
          + \tfrac{1}{30}\Delta R\,\mathrm{id}
                                                      + \tfrac{1}{12}\Omega_{\mu\nu}\Omega^{\mu\nu}
                                                      + \text{(curvature}^2\text{ terms)}
        \right),
      \end{align}
      where $\Omega_{\mu\nu}$ is the curvature of $\nabla^{S,A}$ and the curvature$^2$
      terms are tabulated in~\cite{Vassilevich2003}, Section~4.

      The $\gamma_5$-weighted variant replaces $\operatorname{tr}$ by
      $\operatorname{tr}(\gamma_5\,\cdot\,)$.
      Only terms containing four Clifford matrices survive; the gauge-curvature part
      yields the topological density~\eqref{eq:topological-density}, which controls the
      anomaly constraints~\eqref{eq:anomaly-linear}--\eqref{eq:anomaly-mixed}.

  \end{appendix}

  \newpage
  \bibliographystyle{plain}
  \bibliography{cosmochrony-bibliography,external-refs}

\end{document}
